\documentclass{article}
\usepackage{geometry}
\geometry{margin=1.5cm, vmargin={0pt,1cm}}
\setlength{\topmargin}{-1cm}
\setlength{\headheight}{12.5pt}
\setlength{\paperheight}{29.7cm}
\setlength{\textheight}{25.3cm}

% Chinese support (must load before macros to avoid fontspec conflicts)
\usepackage[UTF8]{ctex}

% useful packages.
\usepackage{fancyhdr}
\usepackage{layout}

% import common macros
% % % % % % % % % % % % % % % % % % % % % % % % % % % % % % % %
% Common Macros for LaTeX
% % % % % % % % % % % % % % % % % % % % % % % % % % % % % % % %

% Useful packages
\usepackage{amsfonts}
\usepackage{amsmath}
\usepackage{amssymb}
\usepackage{amsthm}
\usepackage{enumerate}
\usepackage{graphicx}
\usepackage{multicol}
\usepackage[ruled,linesnumbered]{algorithm2e}

% Math operators and common symbols
\newcommand{\dif}{\mathrm{d}}
\newcommand{\innerProd}[1]{\left\langle #1 \right\rangle}
\newcommand{\difFrac}[2]{\frac{\dif #1}{\dif #2}}
\newcommand{\pdfFrac}[2]{\frac{\partial #1}{\partial #2}}
\newcommand{\T}{\mathrm{T}}
\newcommand{\norm}[1]{\left\| #1 \right\|}
\newcommand{\abs}[1]{\left| #1 \right|}

% Vector and Matrix shortcuts
\newcommand{\vecTwo}[2]{
  \begin{bmatrix}
    #1 \\
    #2
\end{bmatrix}}
\newcommand{\vecThree}[3]{
  \begin{bmatrix}
    #1 \\
    #2 \\
    #3
\end{bmatrix}}
\newcommand{\vecTwoH}[2]{
  \begin{bmatrix}
    #1 & #2
  \end{bmatrix}
}
\newcommand{\vecThreeH}[3]{
  \begin{bmatrix}
    #1 & #2 & #3
  \end{bmatrix}
}

% Common fields
\newcommand{\R}{\mathbb{R}}
\newcommand{\C}{\mathbb{C}}
\newcommand{\Z}{\mathbb{Z}}
\newcommand{\N}{\mathbb{N}}

% Solutions and proofs
\newenvironment{solution}{
\begin{proof}[解]}{
\end{proof}}

% Utility to get environment variables (requires lualatex)
\newcommand{\getEnv}[2]{\directlua{tex.print(os.getenv("#1") or "#2")}}


% Override ctex's Chinese proof label
\renewcommand{\proofname}{Proof}

% Clickable cross-references
\usepackage{xcolor}
\usepackage{hyperref}
\hypersetup{colorlinks=true, linkcolor=blue!60!black, urlcolor=blue!60!black}

\begin{document}

\pagestyle{fancy}
\fancyhead{}
\lhead{Point Set Topology --- Midterm Reference}
\rhead{\today}

\begin{center}
  {\Large\bfseries Point Set Topology --- Midterm Exam Reference Sheet}\\[6pt]
  {\large Table of Contents}
\end{center}
\vspace{6pt}

\tableofcontents
\newpage

%=====================================================================
\section{集合论与公理 (Set Theory and ZFC Axioms)}
%=====================================================================

\subsection{Axiom of Extensionality (外延公理)}\label{def:extensionality}
Two sets are equal if and only if they have the same elements:
\[
  A = B \iff (\forall x)(x \in A \Leftrightarrow x \in B).
\]

\subsection{Axiom of Empty Set (空集公理)}\label{def:empty-set}
There exists a set containing no elements:
\[
  (\exists A)(\forall x)(x \notin A).
\]
This unique set is denoted $\varnothing$ (the empty set).

\subsection{Axiom of Pairing (配对公理)}\label{def:pairing}
Given any two sets $A$ and $B$, there exists a set $\{A, B\}$
containing exactly $A$ and $B$.

\subsection{Axiom of Union (并集公理)}\label{def:union-axiom}
For any set $A$, there exists a set $\bigcup A = \{x \mid (\exists B
\in A)(x \in B)\}$.
Example: if $A = \{\{1,2\}, \{3\}\}$, then $\bigcup A = \{1, 2, 3\}$.

\subsection{Axiom of Power Set (幂集公理)}\label{def:power-set-axiom}
For any set $A$, the collection of all subsets of $A$ forms a set:
\[
  \mathcal{P}(A) = \{B \mid B \subseteq A\}.
\]

\subsection{Axiom Schema of Separation (分离公理/概括公理)}\label{def:separation}
Given a set $A$ and a property $P(x)$, the subset $\{x \in A \mid
P(x)\}$ exists.
\textbf{Key:} must start from an existing set $A$; this avoids
Russell's paradox.

\subsection{Axiom Schema of Replacement (替换公理)}\label{def:replacement}
If $\varphi(x, y)$ defines a ``functional property'' (for each $x \in
A$ there exists a unique $y$), then all such $y$ form a set.

\subsection{Axiom of Infinity (无穷公理)}\label{def:infinity}
There exists an infinite set. Formally: $\exists I$ such that
$\varnothing \in I$ and $(\forall x \in I)(x \cup \{x\} \in I)$.
This constructs $\N$: $0 = \varnothing$, $1 = \{0\}$, $2 = \{0, 1\}$, \ldots

\subsection{Axiom of Regularity / Foundation (正则公理)}\label{def:regularity}
Every nonempty set $A$ has an element disjoint from $A$:
\[
  (\forall A \neq \varnothing)(\exists x \in A)(x \cap A = \varnothing).
\]
\textbf{Corollary:} No set satisfies $A \in A$ (no ``infinite Russian
doll'' self-containment).

\subsection{Axiom of Choice --- AC (选择公理)}\label{def:ac}
If $A$ is a family of nonempty sets, there exists a function $f : A
\to \bigcup A$ such that $(\forall X \in A)(f(X) \in X)$.
One simultaneously chooses one element from each set.

\textbf{Independence:} G\"odel (1938) showed AC is consistent with
ZF; Cohen (1963) showed AC is not provable from ZF. Hence AC is
\emph{independent} of ZF.

\subsection{Russell's Paradox (罗素悖论)}\label{def:russell}
If one replaces Separation with unrestricted comprehension $\{x \mid
P(x)\}$ is always a set, let $A = \{x \mid x \notin x\}$. Then $A \in
A \Leftrightarrow A \notin A$, a contradiction. ZFC avoids this via
the Separation axiom.

\subsection{Forcing (力迫法)}\label{def:forcing}
A technique invented by Cohen (1963) to prove independence results in
set theory. Used to show both $\neg$AC and $\neg$CH are consistent with ZF.

%=====================================================================
\section{可数性与基数 (Countability and Cardinality)}
%=====================================================================

\subsection{Finite Set (有限集)}\label{def:finite}
A set $S$ is \textbf{finite} if there exists a bijection $S \to \{1,
2, \ldots, n\}$ for some $n \in \N$. The number $n$ is the \textbf{size} $|S|$.

\subsection{Infinite Set (无限集)}\label{def:infinite}
A set is \textbf{infinite} if it is not finite.

\subsection{Countably Infinite (可数无限)}\label{def:countably-infinite}
A set $S$ is \textbf{countably infinite} if there exists a bijection
$S \to \N$. Intuitively: the elements can be listed as $s_0, s_1, s_2, \ldots$

\subsection{Countable (可数)}\label{def:countable}
A set is \textbf{countable} if it is either finite or countably infinite.

\subsection{Uncountable (不可数)}\label{def:uncountable}
A set is \textbf{uncountable} if it is not countable.

\textbf{Key equivalence (for nonempty $B$):}
\[
  B \text{ is countable} \iff \exists\text{ surjection } \N
  \twoheadrightarrow B \iff \exists\text{ injection } B \hookrightarrow \N.
\]

\begin{proof}
  ($B$ countable $\Rightarrow$ surjection $\N \to B$): If $B$ is
  finite $\{b_1, \ldots, b_n\}$,
  define $f : \N \to B$ by $f(k) = b_{((k-1) \bmod n) + 1}$ (cycling).
  If $B$ is countably infinite, there is a bijection $\N \to B$,
  which is a surjection.

  (Surjection $\N \to B$ $\Rightarrow$ injection $B \hookrightarrow \N$):
  Given surjection $g : \N \to B$, for each $b \in B$ pick $n_b \in
  g^{-1}(\{b\})$
  (by AC). Then $b \mapsto n_b$ is injective: if $n_{b_1} = n_{b_2}$,
  then $g(n_{b_1}) = b_1 = b_2$.

  (Injection $B \hookrightarrow \N$ $\Rightarrow$ $B$ countable):
  Let $i : B \hookrightarrow \N$ be injective. Then $B \cong i(B) \subseteq \N$.
  Every subset of $\N$ is countable (well-order $i(B)$ by the natural
    order of $\N$;
  if infinite, enumerate in increasing order to get a bijection with $\N$).
\end{proof}

\subsection{Cantor Diagonal Argument (Cantor 对角线论证)}\label{def:cantor-diag}
\begin{theorem}
  $\{0, 1\}^{\N}$ (the set of all infinite binary sequences) is uncountable.
\end{theorem}
\begin{proof}
  Assume countable. List all sequences $x_1, x_2, x_3, \ldots$ where
  $x_i = (x_{i1}, x_{i2}, \ldots)$.
  Define $y = (y_1, y_2, \ldots)$ by $y_j \neq x_{jj}$ for all $j$.
  Then $y \neq x_i$ for all $i$ (differs at position $i$), contradiction.
\end{proof}

\subsection{Cantor's Theorem (Cantor 定理)}\label{def:cantor-thm}
\begin{theorem}
  For any set $A$, there is no surjection $A \twoheadrightarrow
  \mathcal{P}(A)$. Hence $|A| < |\mathcal{P}(A)|$.
\end{theorem}
\begin{proof}
  Suppose $g : A \to \mathcal{P}(A)$ is a surjection. Define $B = \{a
  \in A : a \notin g(a)\}$.
  Since $g$ is surjective, $\exists b \in A$ with $g(b) = B$.
  If $b \in B = g(b)$, then by definition of $B$, $b \notin g(b) =
  B$, contradiction.
  If $b \notin B = g(b)$, then by definition of $B$, $b \in B$, contradiction.
\end{proof}

\subsection{Cardinal Number and Cardinality (基数与势)}\label{def:cardinality}
Two sets $A, B$ have the \textbf{same cardinality} (same cardinal
number), written $|A| = |B|$, if there exists a bijection $A \to B$.

$|A| < |B|$ means: there exists an injection $A \hookrightarrow B$
but no bijection $A \to B$.

\subsection{$\aleph_0$ and $\aleph_1$}\label{def:aleph}
$\aleph_0 = |\N|$ (cardinality of countably infinite sets).
$\aleph_1 = |\R| = |\mathcal{P}(\N)|$ (cardinality of the continuum).

\subsection{Continuum Hypothesis --- CH (连续统假设)}\label{def:ch}
There is no set $A$ with $\aleph_0 < |A| < \aleph_1$.
CH is \textbf{independent} of ZFC (G\"odel 1940: consistent; Cohen
1963: not provable).

%=====================================================================
\section{拓扑空间 (Topological Spaces)}
%=====================================================================

\subsection{Topological Space and Topology (拓扑空间与拓扑)}\label{def:topology}
\begin{definition}
  A \textbf{topology} on a set $X$ is a collection $\mathcal{T}
  \subseteq \mathcal{P}(X)$ satisfying:
  \begin{enumerate}
    \item $\varnothing \in \mathcal{T}$ and $X \in \mathcal{T}$.
    \item Arbitrary unions: if $\{U_i\}_{i \in I} \subseteq
      \mathcal{T}$, then $\bigcup_{i \in I} U_i \in \mathcal{T}$.
    \item Finite intersections: if $U_1, \ldots, U_k \in
      \mathcal{T}$, then $\bigcap_{i=1}^k U_i \in \mathcal{T}$.
  \end{enumerate}
  The pair $(X, \mathcal{T})$ is a \textbf{topological space}.
\end{definition}

\subsection{Open Set (开集)}\label{def:open}
A subset $U \subseteq X$ is \textbf{open} (or $\mathcal{T}$-open) if
$U \in \mathcal{T}$.
``Open'' is relative to the topology, not absolute.

\subsection{Closed Set (闭集)}\label{def:closed}
A subset $C \subseteq X$ is \textbf{closed} if $X \setminus C$ is open.

\textbf{Properties of closed sets:}
\begin{itemize}
  \item Arbitrary intersections of closed sets are closed.
  \item Finite unions of closed sets are closed.
\end{itemize}
\begin{proof}
  Let $\{C_\alpha\}$ be closed sets. Then $X \setminus \bigcap_\alpha
  C_\alpha = \bigcup_\alpha (X \setminus C_\alpha)$
  is a union of open sets, hence open. So $\bigcap_\alpha C_\alpha$ is closed.
  Similarly, $X \setminus (C_1 \cup \cdots \cup C_k) = (X \setminus
  C_1) \cap \cdots \cap (X \setminus C_k)$
  is a finite intersection of open sets, hence open.
\end{proof}

\subsection{Clopen Set (既开又闭集)}\label{def:clopen}
A set that is simultaneously open and closed.
$\varnothing$ and $X$ are always clopen. The existence of
\emph{other} clopen sets relates to disconnectedness.

\subsection{Trivial / Indiscrete Topology (平凡拓扑/密着拓扑)}\label{def:trivial}
$\mathcal{T} = \{\varnothing, X\}$. The coarsest possible topology.

\subsection{Discrete Topology (离散拓扑)}\label{def:discrete-topo}
$\mathcal{T} = \mathcal{P}(X)$. Every subset is open. The finest
possible topology.

\subsection{Cofinite Topology (余有限拓扑)}\label{def:cofinite}
$\mathcal{T} = \{\varnothing\} \cup \{X \setminus F : F \text{ is
finite}\}$. The open sets are $\varnothing$ and complements of finite sets.

\begin{proof}[Verification that this is a topology]
  $\varnothing, X \in \mathcal{T}$ trivially.
  Arbitrary unions: let $\{X \setminus F_\alpha\}$ be open sets.
  Then $\bigcup_\alpha (X \setminus F_\alpha) = X \setminus
  \bigcap_\alpha F_\alpha$.
  Since $\bigcap_\alpha F_\alpha \subseteq F_{\alpha_0}$ is finite,
  the union is open.
  Finite intersections: $(X \setminus F_1) \cap (X \setminus F_2) = X
  \setminus (F_1 \cup F_2)$,
  and $F_1 \cup F_2$ is finite.
\end{proof}

\subsection{Semicontinuous Topology (半无穷拓扑)}\label{def:semicontinuous}
On $\R$: $\mathcal{T}$ consists of $\varnothing$, $\R$, and all sets
of the form $(a, +\infty)$ (and their finite intersections, which are
still of this form).
This is coarser than the standard topology on $\R$.

\subsection{Finer / Coarser (更细/更粗)}\label{def:finer-coarser}
$\mathcal{T}_1$ is \textbf{finer} (larger) than $\mathcal{T}_2$ if
$\mathcal{T}_2 \subseteq \mathcal{T}_1$.
$\mathcal{T}_1$ is \textbf{coarser} (smaller) than $\mathcal{T}_2$ if
$\mathcal{T}_1 \subseteq \mathcal{T}_2$.

\textbf{Fact:} A finer topology has more open sets, hence more
continuous functions. A coarser topology has fewer.

%=====================================================================
\section{基 (Basis)}
%=====================================================================

\subsection{Basis (基)}\label{def:basis}
\begin{definition}
  $\mathcal{B} \subseteq \mathcal{P}(X)$ is a \textbf{basis for a topology} if:
  \begin{enumerate}
    \item (Covering) $\bigcup_{B \in \mathcal{B}} B = X$.
    \item (Intersection) If $x \in B_1 \cap B_2$ with $B_1, B_2 \in
      \mathcal{B}$, then $\exists B_3 \in \mathcal{B}$ with $x \in
      B_3 \subseteq B_1 \cap B_2$.
  \end{enumerate}
\end{definition}
The topology \textbf{generated} by $\mathcal{B}$ is $\mathcal{T} =
\{\bigcup_{i \in I} B_i : B_i \in \mathcal{B}\}$.

\subsection{Subbasis (子基)}\label{def:subbasis}
$\mathcal{S} \subseteq \mathcal{P}(X)$ is a \textbf{subbasis} if the
collection of all finite intersections of elements of $\mathcal{S}$
forms a basis.
The topology generated by $\mathcal{S}$ is the coarsest topology
containing $\mathcal{S}$.

\subsection{Generating a Topology (生成拓扑)}\label{def:generate-topo}
Given a basis $\mathcal{B}$, every open set is a union of basis
elements: $U = \bigcup\{B \in \mathcal{B} : B \subseteq U\}$.

\subsection{Countable Basis / Second-Countable
(可数基/第二可数)}\label{def:second-countable}
A space $(X, \mathcal{T})$ is \textbf{second-countable} if its
topology has a countable basis.
Example: $\R$ with the basis $\{(a, b) : a, b \in \Q\}$ (rational endpoints).

\subsection{First-Countable (第一可数)}\label{def:first-countable}
\begin{definition}
  $(X, \mathcal{T})$ is \textbf{first-countable} if every point
  $x \in X$ has a \emph{countable neighborhood basis} (local basis): a
  countable collection $\{U_n\}_{n \in \N}$ of open neighborhoods of $x$
  such that every open $V \ni x$ contains some $U_n$.
\end{definition}

\textbf{Intuition:} Second-countability is a \emph{global} condition (one
countable basis for the entire topology).  First-countability is a \emph{local}
condition (each point merely needs its own countable family of neighborhoods).

\textbf{Key facts:}
\begin{itemize}
  \item \textbf{Second-countable $\Rightarrow$ first-countable}
    (a global countable basis restricts to a countable local basis at each
    point), but \emph{not} conversely.
  \item Every \textbf{metrizable} space is first-countable:
    $\{B(x, 1/n) : n \in \N\}$ is a countable neighborhood basis at $x$.
  \item First-countability is \textbf{hereditary}: every subspace of a
    first-countable space is first-countable.
  \item In a first-countable space, \textbf{sequences suffice} for topology:
    $x \in \overline{A}$ iff there exists a sequence in $A$ converging to $x$;
    $f$ is continuous iff it preserves sequential convergence.
    (In non-first-countable spaces, nets or filters are needed.)
\end{itemize}

\textbf{Examples showing strictness of the hierarchy:}
\begin{itemize}
  \item $\R$ with the standard topology: second-countable
    (hence first-countable).
  \item An uncountable set with the discrete topology: first-countable
    (each point $\{x\}$ is itself a neighborhood basis), but \emph{not}
    second-countable (would require uncountably many basis elements to cover
    all singletons).
  \item $\R_\ell$ (lower limit topology): first-countable
    ($\{[x, x+1/n)\}$ is a local basis at $x$) and separable, but \emph{not}
    second-countable.
  \item The cocountable topology on an uncountable set: \emph{not}
    first-countable (no point has a countable neighborhood basis), hence not
    second-countable.
\end{itemize}

\textbf{First-countable vs.\ second-countable --- summary:}

\begin{center}
  \begin{tabular}{lcc}
    \textbf{Property} & \textbf{First-countable} & \textbf{Second-countable} \\
    \hline
    Scope & Local (per point) & Global (whole space) \\
    Metrizable $\Rightarrow$ & Yes & No ($\R^\omega$ is metrizable
    but not $2$nd-count.) \\
    Hereditary & Yes & Yes \\
    Products & Preserved (finite/countable) & Preserved (countable) \\
    Sequences determine topology & Yes & Yes (strictly stronger) \\
  \end{tabular}
\end{center}

\subsection{Open Rectangle / Open Ball (开矩形/开球)}\label{def:open-rect}
In $\R^n$ with Euclidean metric: the \textbf{open ball} $B(x,
\varepsilon) = \{y : d(x, y) < \varepsilon\}$.
The \textbf{open rectangle} $(a_1, b_1) \times \cdots \times (a_n,
b_n)$ is a basis element for the product topology.

%=====================================================================
\section{分离公理 (Separation Axioms: $T_0$--$T_4$)}
%=====================================================================

The separation axioms are a hierarchy of increasingly strong conditions on a
topological space, measuring the degree to which distinct points (or points and
closed sets, or pairs of closed sets) can be ``separated'' by open sets.  These
conditions are important in their own right and are the central tool for
characterizing metrizability.

\subsection{Overview: $T_0 \Rightarrow T_1 \Rightarrow T_2
\Rightarrow T_3 \Rightarrow T_{3\frac{1}{2}} \Rightarrow T_4$}

\begin{center}
  \begin{tabular}{lll}
    \textbf{Name} & \textbf{Symbol} & \textbf{Requirement} \\
    \hline
    Kolmogorov   & $T_0$     & Distinct points are topologically
    distinguishable \\
    Fr\'echet    & $T_1$     & Singletons are closed \\
    Hausdorff    & $T_2$     & Distinct points have disjoint open
    neighborhoods \\
    Regular      & $T_3$     & Points and disjoint closed sets
    separated by open sets \\
    Tychonoff    & $T_{3\frac{1}{2}}$ & Points and disjoint closed
    sets separated by continuous functions \\
    Normal       & $T_4$     & Disjoint closed sets separated by open sets \\
  \end{tabular}
\end{center}

\textbf{Convention:} $T_i$ denotes both the axiom and the space satisfying it.
``$T_3$'' = regular + $T_1$; ``$T_4$'' = normal + $T_1$. Some authors write
``regular'' for $T_3$ without $T_1$; we follow the convention that
$T_3$ \emph{includes} $T_1$.

\subsection{$T_0$ Space --- Kolmogorov Space (Kolmogorov 空间)}\label{def:T0}
\begin{definition}
  $(X, \mathcal{T})$ is $T_0$ if for any two distinct points $x \ne y$,
  there exists an open set $U$ containing exactly one of them.
\end{definition}

\textbf{Intuition:} Topology can ``distinguish'' any two points --- you
never have two points with \emph{exactly} the same set of open neighborhoods.

\textbf{Non-example:} The trivial (indiscrete) topology on $\{a, b\}$ is
\emph{not} $T_0$: the only open sets are $\varnothing$ and $\{a, b\}$,
so $a$ and $b$ have identical neighborhoods.

\textbf{Example:} The Sierpi\'nski space $\{0, 1\}$ with topology
$\{\varnothing, \{1\}, \{0,1\}\}$ is $T_0$ but not $T_1$: $\{0\}$ is not closed.

\subsection{$T_1$ Space --- Fr\'echet Space (Fr\'echet 空间)}\label{def:T1}
\begin{definition}
  $(X, \mathcal{T})$ is $T_1$ if for any two distinct points $x \ne y$,
  there exist open sets $U \ni x$, $V \ni y$ with $y \notin U$ and $x \notin V$.
\end{definition}

\begin{proposition}
  The following are equivalent:
  \begin{enumerate}
    \item $X$ is $T_1$.
    \item Every singleton $\{x\}$ is closed.
    \item Every finite subset of $X$ is closed.
  \end{enumerate}
\end{proposition}
\begin{proof}
  ($1 \Rightarrow 2$): For $x \in X$ and $y \ne x$, by $T_1$ there exists
  open $V_y \ni y$ with $x \notin V_y$. Then $X \setminus \{x\} =
  \bigcup_{y \ne x} V_y$ is open, so $\{x\}$ is closed.

  ($2 \Rightarrow 3$): Finite unions of closed sets are closed.

  ($3 \Rightarrow 1$): Given $x \ne y$, the set $\{y\}$ is closed, so
  $U = X \setminus \{y\}$ is open, $x \in U$, $y \notin U$.
  Symmetrically for $V$.
\end{proof}

\textbf{Key consequences of $T_1$:}
\begin{itemize}
  \item A sequence can have \emph{at most one} limit (limits need not exist,
    but they are unique when they do).
  \item $T_1$ is \textbf{hereditary}: every subspace of a $T_1$ space is $T_1$.
  \item Products of $T_1$ spaces are $T_1$.
\end{itemize}

\textbf{Example:} Any metric space is $T_1$ (indeed $T_2$, $T_4$).
The cofinite topology on an infinite set is $T_1$ but \emph{not} $T_2$.

\subsection{$T_2$ Space --- Hausdorff Space (Hausdorff 空间)}\label{def:T2}
\begin{definition}
  $(X, \mathcal{T})$ is \textbf{Hausdorff} ($T_2$) if for any two distinct
  points $x \ne y$, there exist disjoint open sets $U, V$ with
  $x \in U$, $y \in V$, i.e.\ $U \cap V = \varnothing$.
\end{definition}

\begin{proposition}
  In a Hausdorff space:
  \begin{enumerate}
    \item Every convergent sequence has a \textbf{unique} limit.
    \item Every compact subset is \textbf{closed}.
    \item The diagonal $\Delta = \{(x,x) : x \in X\}$ is closed in $X \times X$.
  \end{enumerate}
\end{proposition}

\textbf{Proof sketch (compact $\Rightarrow$ closed):} Let $K \subseteq Y$ be
compact, $Y$ Hausdorff, and $y \notin K$. For each $x \in K$, pick disjoint
open $U_x \ni x$, $V_x \ni y$. The $U_x$ cover $K$, so finitely many
$U_{x_1}, \ldots, U_{x_n}$ suffice. Then $V = V_{x_1} \cap \cdots \cap V_{x_n}$
is an open neighborhood of $y$ disjoint from $K$. So $Y \setminus K$ is open.

\textbf{Key facts:}
\begin{itemize}
  \item Hausdorff is \textbf{hereditary}: every subspace of a
    Hausdorff space is Hausdorff.
  \item Products of Hausdorff spaces are Hausdorff.
  \item A continuous bijection from a \textbf{compact} space to a
    \textbf{Hausdorff} space is a homeomorphism
    (this is one of the most-used tools in point-set topology).
\end{itemize}

\textbf{Example:} $\R^n$, any metric space, any order topology.
\textbf{Non-example:} Cofinite topology on infinite set; Zariski
topology on $\C^n$.

\subsection{$T_3$ Space --- Regular Space (正则空间)}\label{def:T3}
\begin{definition}
  $(X, \mathcal{T})$ is \textbf{regular} ($T_3$) if it is $T_1$ and for
  every point $x$ and closed set $F$ with $x \notin F$, there exist
  disjoint open sets $U \ni x$ and $V \supseteq F$.
\end{definition}

\textbf{Equivalent reformulation:} $X$ is regular iff each point $x$ has
a neighborhood basis consisting of \emph{closed neighborhoods}
$\bar{U}$ --- i.e., for every open $U \ni x$, there exists open $V$ with
$x \in V \subseteq \bar{V} \subseteq U$.

\textbf{Key facts:}
\begin{itemize}
  \item Regularity is hereditary: every subspace of a regular space is regular.
  \item Products of regular spaces are regular.
  \item Every metric space is regular (indeed normal).
  \item \textbf{Compact + Hausdorff $\Rightarrow$ regular} (and normal).
\end{itemize}

\subsection{$T_{3\frac{1}{2}}$ Space --- Tychonoff / Completely Regular Space
(完全正则空间)}\label{def:T3h}
\begin{definition}
  $(X, \mathcal{T})$ is \textbf{completely regular} ($T_{3\frac{1}{2}}$)
  if it is $T_1$ and for every point $x$ and closed set $F$ with $x \notin F$,
  there exists a continuous function $f : X \to [0,1]$ with $f(x) = 0$
  and $f(F) = \{1\}$.
\end{definition}

This lies strictly between $T_3$ and $T_4$:
\[ T_4 \;\Rightarrow\; T_{3\frac{1}{2}} \;\Rightarrow\; T_3. \]

\textbf{Key fact (Urysohn's Lemma $\Rightarrow$ $T_4$ is $T_{3\frac{1}{2}}$):}
Every normal space is completely regular (by Urysohn's Lemma).

\textbf{Why $T_{3\frac{1}{2}}$ matters:} A space embeds into a (Tychonoff)
cube $[0,1]^J$ for some index set $J$ if and only if it is $T_{3\frac{1}{2}}$.
This is the content of the \emph{Tychonoff Embedding Theorem}.

\subsection{$T_4$ Space --- Normal Space (正规空间)}\label{def:T4}
\begin{definition}
  $(X, \mathcal{T})$ is \textbf{normal} ($T_4$) if it is $T_1$ and for
  every pair of disjoint closed sets $A, B$, there exist disjoint open
  sets $U \supseteq A$, $V \supseteq B$.
\end{definition}

\begin{theorem}[Urysohn's Lemma]
  $X$ is normal if and only if for every pair of disjoint closed sets
  $A, B \subseteq X$, there exists a continuous function
  $f : X \to [0,1]$ with $f(A) = \{0\}$ and $f(B) = \{1\}$.
\end{theorem}

This is one of the most important results in point-set topology. It connects
the purely topological notion of ``separating closed sets by open sets'' with
the analytic notion of ``separating by continuous functions.''

\begin{theorem}[Tietze Extension Theorem]
  $X$ is normal if and only if for every closed $A \subseteq X$ and every
  continuous $f : A \to \R$, there exists a continuous extension
  $F : X \to \R$ with $F|_A = f$.
\end{theorem}

\textbf{Key facts about normality:}
\begin{itemize}
  \item \textbf{Compact + Hausdorff $\Rightarrow$ normal.}
    (Proof: Let $A, B$ be disjoint closed sets in a compact Hausdorff space.
      $A$ is compact (closed in compact). Use the Hausdorff property to
    separate each $a \in A$ from $B$, extract a finite subcover.)
  \item \textbf{Metrizable $\Rightarrow$ normal.}
    (Proof: see \S\ref{def:metrizable}, the function
    $f(x) = d(x,A)/(d(x,A)+d(x,B))$ does the job.)
  \item Normality is \textbf{NOT hereditary} in general
    (a subspace of a normal space need not be normal).
  \item Normality is \textbf{NOT preserved by products} in general
    ($\R^\omega$ in the product topology is normal, but
    Sorgenfrey plane $\mathbb{S} \times \mathbb{S}$ is not).
\end{itemize}

\textbf{Famous non-normal space:} The \emph{Sorgenfrey line}
$\mathbb{S}$ (topology generated by $[a, b)$) is $T_3$ (regular) but
the product $\mathbb{S} \times \mathbb{S}$ is \emph{not} $T_4$ (not normal).
Also, $\R^\omega$ in the box topology is not normal
(see \S\ref{def:metrizable}).

\textbf{Jones' Lemma:} If $X$ is normal and contains a dense set $D$ and
a closed discrete set $S$ with $|S| \ge 2^{|D|}$, then $X$ is not normal.
(Used to show certain spaces are not $T_4$.)

\subsection{Separation Axioms and Metrizability --- Summary
(分离公理与可度量化的关系)}\label{def:sep-metrizable}

The hierarchy of separation axioms gives both \textbf{necessary conditions}
and (combined with countability) \textbf{sufficient conditions} for
metrizability:

\begin{center}
  \begin{tabular}{ll}
    \textbf{Direction} & \textbf{Statement} \\
    \hline
    Necessary & Metrizable $\Rightarrow$ $T_4$ (normal) \\
    Necessary & Metrizable $\Rightarrow$ first-countable \\
    Necessary & Metrizable $\Rightarrow$ $T_2$ (Hausdorff) \\
    \hline
    Sufficient & $T_3$ + second-countable $\Rightarrow$ metrizable (Urysohn) \\
    Sufficient & Compact + metrizable $\Rightarrow$ second-countable \\
  \end{tabular}
\end{center}

\begin{theorem}[Urysohn Metrization Theorem]
  \label{thm:urysohn-metrization}
  If $(X, \mathcal{T})$ is $T_3$ (regular $+$ $T_1$) and second-countable,
  then $X$ is metrizable.
\end{theorem}

\textbf{Proof idea:} Use second-countability to embed $X$ into
the Hilbert cube $[0,1]^\omega$ (a metrizable space). The regularity
condition ensures we can separate points from closed sets using
Urysohn functions, constructing an embedding.

\textbf{Converse direction:}
\begin{itemize}
  \item Metrizable + separable $\Rightarrow$ second-countable
    (the countable dense set gives a countable basis via $B(d, 1/n)$).
  \item Metrizable + compact $\Rightarrow$ second-countable
    (compact $\Rightarrow$ totally bounded $\Rightarrow$ separable,
    then apply above).
\end{itemize}

So for \emph{compact} spaces: \textbf{metrizable $\iff$ second-countable
$\iff$ separable.}

\subsection{Second-Countability Revisited
(第二可数再探)}\label{def:second-countable-revisit}

Recall (\S\ref{def:second-countable}): $X$ is \textbf{second-countable}
if $\mathcal{T}$ has a countable basis.

\textbf{Equivalent / related conditions for second-countability:}
\begin{itemize}
  \item $X$ has a countable basis $\iff$ $X$ is \textbf{separable}
    (has a countable dense subset) \textbf{and} metrizable.
    (In general, second-countable $\Rightarrow$ separable, but not conversely
    unless metrizable.)
  \item $X$ is second-countable $\Rightarrow$ $X$ is \textbf{Lindel\"of}
    (every open cover has a countable subcover).
  \item $X$ is second-countable $\Rightarrow$ $X$ is \textbf{first-countable}
    (each point has a countable neighborhood basis).
\end{itemize}

\textbf{Counterexamples to show strictness:}
\begin{itemize}
  \item $\R$ with the standard topology: second-countable
    (basis $\{(a,b) : a, b \in \Q\}$).
  \item $\R_\ell$ (lower limit topology, basis $\{[a,b)\}$):
    separable ($\Q$ is dense) and first-countable, but \emph{not}
    second-countable.
  \item An uncountable set with the discrete topology:
    first-countable (trivially) but not second-countable.
\end{itemize}

\textbf{Preservation properties:}
\begin{itemize}
  \item Second-countability is \textbf{hereditary}: every subspace of a
    second-countable space is second-countable.
  \item Countable \textbf{products} of second-countable spaces are
    second-countable.
  \item Second-countability is a \textbf{topological invariant}
    (preserved by homeomorphisms).
\end{itemize}

%=====================================================================
\section{度量空间 (Metric Spaces)}
%=====================================================================

\subsection{Metric Space and Metric (度量空间与度量)}\label{def:metric}
\begin{definition}
  A \textbf{metric space} $(S, d)$ consists of a set $S$ and a
  function $d : S \times S \to \R$ satisfying:
  \begin{enumerate}
    \item $d(x, y) \ge 0$, with $d(x, y) = 0 \iff x = y$.
    \item $d(x, y) = d(y, x)$ (symmetry).
    \item $d(x, z) \le d(x, y) + d(y, z)$ (triangle inequality).
  \end{enumerate}
\end{definition}

\subsection{Discrete Metric (离散度量)}\label{def:discrete-metric}
$d(x, y) =
\begin{cases} 0 & x = y \\ 1 & x \ne y
\end{cases}$. Any two distinct points are at distance 1. This induces
the discrete topology.

\subsection{Open Ball and Closed Ball (开球与闭球)}\label{def:balls}
\textbf{Open ball:} $B(x, \varepsilon) = \{y \in S : d(x, y) < \varepsilon\}$.
\textbf{Closed ball:} $\bar{B}(x, \varepsilon) = \{y \in S : d(x, y)
\le \varepsilon\}$.

\subsection{Bounded Set (有界集)}\label{def:bounded}
$A \subseteq (S, d)$ is \textbf{bounded} if $\exists M > 0$ such that
$d(x, y) \le M$ for all $x, y \in A$.
Equivalently: $\text{diam}(A) = \sup\{d(x, y) : x, y \in A\} < \infty$.

\subsection{Bounded Metric (有界度量)}\label{def:bounded-metric}
\begin{proposition}
  For any metric $d$, the function $d^*(x, y) = \min\{d(x, y), 1\}$
  is a metric that induces the same topology as $d$.
\end{proposition}
\textbf{Warning:} There need not exist a uniform $A > 0$ with $A
\cdot d \le d^*$ (e.g., $X = \R$).

\begin{proof}[Proof that $d^*$ is a metric]
  Non-negativity and symmetry are immediate. For the triangle inequality:
  $d^*(x,z) = \min\{d(x,z),1\} \le \min\{d(x,y)+d(y,z), 1\}$.
  If $d(x,y) \ge 1$ or $d(y,z) \ge 1$, then $d^*(x,y)+d^*(y,z) \ge 1
  \ge d^*(x,z)$.
  If both $< 1$, then $d^*(x,z) \le d(x,z) \le d(x,y)+d(y,z) =
  d^*(x,y)+d^*(y,z)$.
\end{proof}

\begin{proof}[Proof that $d^*$ induces the same topology]
  For $0 < r < 1$: $B_{d^*}(x,r) = \{y : d^*(x,y) < r\} = \{y :
  d(x,y) < r\} = B_d(x,r)$,
  so $d^*$-balls with radius $< 1$ are exactly $d$-balls.
  Any $d^*$-ball $B_{d^*}(x,r)$ with $r \ge 1$ equals all of $X$
  (since $d^* \le 1 < r$),
  which is open in $\mathcal{T}_d$. Conversely, any $d$-ball
  $B_d(x,r)$ equals $B_{d^*}(x,r)$ if $r < 1$,
  or $B_{d^*}(x,1/2) \subseteq B_d(x,r)$ if $r \ge 1$.
\end{proof}

\subsection{Equivalent Metrics (等价度量)}\label{def:equiv-metrics}
Two metrics $d_1, d_2$ on $X$ are \textbf{equivalent} (topologically)
if they induce the same topology.
\begin{theorem}[Sufficient condition]
  If $\exists c_1, c_2 > 0$ such that $c_1 d_1(x,y) \le d_2(x,y) \le
  c_2 d_1(x,y)$ for all $x, y$, then $d_1$ and $d_2$ are equivalent.
\end{theorem}
\begin{proof}
  For any $x$ and $r > 0$:
  $B_{d_2}(x, r) \supseteq B_{d_1}(x, r/c_2)$
  (since $d_1(x,y) < r/c_2$ implies $d_2(x,y) \le c_2 d_1(x,y) < r$).
  Similarly, $B_{d_1}(x, r) \supseteq B_{d_2}(x, c_1 r)$
  (since $d_2(x,y) < c_1 r$ implies $c_1 d_1(x,y) \le d_2(x,y) < c_1
  r$, so $d_1(x,y) < r$).
  So every $d_2$-ball contains a $d_1$-ball and vice versa, hence the
  topologies coincide.
\end{proof}

\subsection{Metrizable Topology (可度量化的拓扑)}\label{def:metrizable}
\begin{definition}
  A topological space $(X, \mathcal{T})$ is \textbf{metrizable} if
  there exists a metric $d$ on $X$ such that $\mathcal{T} = \mathcal{T}_d$.
\end{definition}

\textbf{Key fact:} ``Metrizable'' is a \emph{topological property}:
if $X \cong Y$ and $X$ is metrizable, then $Y$ is metrizable
(pull back the metric along the homeomorphism).

\subsubsection*{Necessary Conditions for Metrizability}

\begin{proposition}
  Every metrizable space is Hausdorff ($T_2$).
\end{proposition}
\begin{proof}
  Let $d$ be a metric inducing the topology.
  For distinct $x, y$: set $r = d(x,y)/2 > 0$.
  Then $B(x,r)$ and $B(y,r)$ are disjoint open neighborhoods
  (if $z \in B(x,r) \cap B(y,r)$, then $d(x,y) \le d(x,z)+d(z,y) < 2r
  = d(x,y)$, contradiction).
\end{proof}

\begin{proposition}
  Every metrizable space is first-countable (each point has a countable
  neighborhood basis).
\end{proposition}
\begin{proof}
  For $x \in X$, the collection $\{B(x, 1/n) : n \in \N\}$ is a countable
  neighborhood basis: any open $U \ni x$ contains $B(x, \varepsilon)$ for some
  $\varepsilon > 0$, and $B(x, 1/n) \subseteq B(x, \varepsilon)$ for
  $n > 1/\varepsilon$.
\end{proof}

\begin{proposition}
  Every metrizable space is normal ($T_4$): disjoint closed sets can be
  separated by disjoint open sets.
\end{proposition}
\begin{proof}
  Let $A, B$ be disjoint closed sets in a metric space $(X,d)$.
  Define $f(x) = \frac{d(x,A)}{d(x,A)+d(x,B)}$ where $d(x,S) =
  \inf_{s \in S} d(x,s)$.
  The denominator is always positive (if $d(x,A)=0$ then $x \in \bar{A}=A$,
  so $x \notin B$, hence $d(x,B)>0$).
  Then $f$ is continuous, $f(A)=\{0\}$, $f(B)=\{1\}$.
  Set $U = f^{-1}([0,1/2))$, $V = f^{-1}((1/2,1])$: these are disjoint open
  sets separating $A$ and $B$.
\end{proof}

\subsubsection*{Sufficient Condition: Urysohn Metrization Theorem}

\begin{theorem}[Urysohn]
  If $(X, \mathcal{T})$ is second-countable and $T_3$ (regular + $T_1$),
  then $X$ is metrizable.
\end{theorem}
(Proof is beyond the scope of this reference. The idea: embed $X$ into the
  Hilbert cube $[0,1]^\omega$ using a countable family of Urysohn functions,
then pull back the metric.)

\textbf{Useful corollaries:}
\begin{itemize}
  \item Compact + Hausdorff + second-countable $\Rightarrow$ metrizable.
  \item A compact metrizable space is second-countable
    (for each $n$, $\{B(x, 1/n) : x \in X\}$ has a finite subcover).
\end{itemize}

\subsubsection*{Proof Techniques}

\textbf{To prove a space IS metrizable:}
\begin{enumerate}
  \item Construct an explicit metric $d$ and show $\mathcal{T} = \mathcal{T}_d$
    (ball comparison: every $d$-ball contains a $\mathcal{T}$-open
    set and vice versa).
    \textbf{Example:} $d_\omega(x,y) = \sup_i \frac{d^*(x_i,y_i)}{i+1}$
    induces the product topology on $\R^\omega$ (see HW5 Q3).
  \item Apply Urysohn's theorem: show second-countable + regular + $T_1$.
  \item Use inheritance: subspaces of metrizable spaces are metrizable
    (restrict the metric); countable products of metrizable spaces
    are metrizable.
\end{enumerate}

\textbf{To prove a space is NOT metrizable}, show it fails a
necessary condition:
\begin{enumerate}
  \item \textbf{Not Hausdorff:} the trivial topology on $\{a,b\}$ is
    not metrizable.
  \item \textbf{Not first-countable:} the box topology on $\R^\omega$ is not
    first-countable (diagonal argument: given any countable neighborhood basis
      $\{U_n\}$ of the origin, each $U_n$ restricts coordinates to
      $(-\varepsilon_{n,i}, \varepsilon_{n,i})$ for each $i$; construct the box
      $V = \prod_i (-\varepsilon_{i,i}/2, \varepsilon_{i,i}/2)$; then
      $V$ is open but
    no $U_n \subseteq V$, since $U_n$ has coordinate $i=n$ wider than $V$).
    Hence not metrizable.
  \item \textbf{Not normal:} the Sorgenfrey line (lower limit topology) is
    $T_3$ but not $T_4$, hence not metrizable.
\end{enumerate}

\subsubsection*{Quick Reference: Metrizability of Common Spaces}

\begin{center}
  \begin{tabular}{@{}lll@{}}
    \hline
    \textbf{Space} & \textbf{Metrizable?} & \textbf{Reason / Metric} \\
    \hline
    $\R^n$, standard topology & Yes & $d_2(x,y) = \|x-y\|$ \\
    Discrete topology on any set & Yes & Discrete metric \\
    Trivial topology, $|X| \ge 2$ & No & Not Hausdorff \\
    Cofinite topology on infinite set & No & Not Hausdorff \\
    $(0,1) \subset \R$ & Yes & Subspace of metrizable \\
    $\R^\omega$, product topology & Yes & $d_\omega = \sup_i d^*_i/(i+1)$ \\
    $\R^\omega$, box topology & No & Not first-countable \\
    $S^n$ (sphere) & Yes & Subspace of $\R^{n+1}$ \\
    \hline
  \end{tabular}
\end{center}

\subsection{Euclidean Metric (欧几里得度量)}\label{def:euclidean}
On $\R^n$: $d_2(x, y) = \sqrt{\sum_{i=1}^n (x_i - y_i)^2}$.

\subsection{Taxicab Metric $\ell^1$ (曼哈顿度量)}\label{def:taxicab}
On $\R^n$: $d_1(x, y) = \sum_{i=1}^n |x_i - y_i|$.

\subsection{Supremum Metric $\ell^\infty$ (最大值度量)}\label{def:supremum}
On $\R^n$: $d_\infty(x, y) = \max_{1 \le i \le n} |x_i - y_i|$.

\textbf{Relations on $\R^n$:}
\[
  d_\infty \le d_2 \le \sqrt{n}\, d_\infty, \qquad d_\infty \le d_1
  \le n\, d_\infty.
\]
All $l^p$ metrics are equivalent in finite dimensions.

\subsection{Norm (范数)}\label{def:norm}
A \textbf{norm} on a vector space $V$ is $\|\cdot\| : V \to [0, \infty)$ with:
$\|x\| = 0 \iff x = 0$, $\|\alpha x\| = |\alpha|\|x\|$, $\|x + y\|
\le \|x\| + \|y\|$.
Every norm induces a metric via $d(x, y) = \|x - y\|$.

%=====================================================================
\section{子空间 (Subspace Topology)}
%=====================================================================

\subsection{Subspace Topology (子空间拓扑)}\label{def:subspace-topo}
\begin{definition}
  If $(X, \mathcal{T})$ is a topological space and $Y \subseteq X$,
  the \textbf{subspace topology} on $Y$ is
  \[
    \mathcal{T}_Y = \{Y \cap U : U \in \mathcal{T}\}.
  \]
\end{definition}
\textbf{Verification} (check four axioms): uses distributivity of
$\cap$ over $\cup$.

\subsection{Subspace (子空间)}\label{def:subspace}
A subset $Y \subseteq X$ endowed with the subspace topology. Open
sets in $Y$ look like $Y \cap U$ where $U$ is open in $X$.

\textbf{Example:} $Y = [1, 2] \subset \R$. The sets $[1, c)$, $(c,
d)$, $(c, 2]$ are all open in $Y$.

\subsection{Hereditary Property (遗传性质)}\label{def:hereditary}
A property $P$ is \textbf{hereditary} if $(X, \mathcal{T})$ has $P$
implies every subspace has $P$.
Examples of hereditary properties: $T_0, T_1, T_2$ (Hausdorff),
second-countability.
Non-example: compactness is not hereditary (only to closed subspaces).

%=====================================================================
\section{闭包、内部与边界 (Closure, Interior, Boundary)}
%=====================================================================

\subsection{Closure (闭包)}\label{def:closure}
\begin{definition}
  The \textbf{closure} of $A \subseteq X$ is $\bar{A} = \bigcap\{C
  \supseteq A : C \text{ is closed}\}$, i.e., the smallest closed set
  containing $A$.
\end{definition}
\textbf{Characterization:} $x \in \bar{A} \iff$ every open
neighborhood of $x$ meets $A$.

\begin{proof}
  ($\Rightarrow$) Suppose $x \in \bar{A}$ and $U$ is an open
  neighborhood of $x$ with $U \cap A = \varnothing$.
  Then $X \setminus U$ is closed and $A \subseteq X \setminus U$.
  Since $\bar{A}$ is the smallest closed superset of $A$, $\bar{A}
  \subseteq X \setminus U$,
  so $x \in X \setminus U$, contradicting $x \in U$.

  ($\Leftarrow$) Suppose every open neighborhood of $x$ meets $A$.
  If $x \notin \bar{A}$, then $X \setminus \bar{A}$ is an open
  neighborhood of $x$
  disjoint from $A$ (since $A \subseteq \bar{A}$), contradiction.
\end{proof}
\textbf{Properties:} $\overline{A \cup B} = \bar{A} \cup \bar{B}$; $A
\subseteq B \Rightarrow \bar{A} \subseteq \bar{B}$; $\bar{\bar{A}} = \bar{A}$.

\begin{proof}[Proof of closure properties]
  \textbf{$\overline{A \cup B} = \bar{A} \cup \bar{B}$:}
  ($\subseteq$) $A \cup B \subseteq \bar{A} \cup \bar{B}$, which is
  closed (finite union of closed).
  Since $\overline{A \cup B}$ is the smallest closed superset,
  $\overline{A \cup B} \subseteq \bar{A} \cup \bar{B}$.
  ($\supseteq$) $A \subseteq A \cup B \subseteq \overline{A \cup B}$,
  so $\bar{A} \subseteq \overline{A \cup B}$.
  Similarly $\bar{B} \subseteq \overline{A \cup B}$, so $\bar{A} \cup
  \bar{B} \subseteq \overline{A \cup B}$.

  \textbf{Monotonicity:} If $A \subseteq B$, then $A \subseteq B
  \subseteq \bar{B}$ (closed).
  Since $\bar{A}$ is the smallest closed superset of $A$, $\bar{A}
  \subseteq \bar{B}$.

  \textbf{Idempotence:} $\bar{A}$ is closed, so $\bar{\bar{A}} =
  \bar{A}$ by definition (closure of a closed set is itself).
\end{proof}

\subsection{Interior (内部)}\label{def:interior}
\begin{definition}
  The \textbf{interior} of $A$ is $A^\circ = \operatorname{int}(A) =
  \bigcup\{U \subseteq A : U \text{ is open}\}$, i.e., the largest
  open set contained in $A$.
\end{definition}
\textbf{Properties:} $\operatorname{int}(A \cap B) =
\operatorname{int}(A) \cap \operatorname{int}(B)$.

\begin{proof}
  ($\subseteq$) $\operatorname{int}(A \cap B) \subseteq A \cap B \subseteq A$,
  and $\operatorname{int}(A \cap B)$ is open, so
  $\operatorname{int}(A \cap B) \subseteq \operatorname{int}(A)$.
  Similarly for $B$, giving $\operatorname{int}(A \cap B) \subseteq
  \operatorname{int}(A) \cap \operatorname{int}(B)$.
  ($\supseteq$) $\operatorname{int}(A) \cap \operatorname{int}(B)$ is
  open (finite intersection of open),
  and $\operatorname{int}(A) \cap \operatorname{int}(B) \subseteq A \cap B$.
  Since $\operatorname{int}(A \cap B)$ is the largest open subset of $A \cap B$,
  $\operatorname{int}(A) \cap \operatorname{int}(B) \subseteq
  \operatorname{int}(A \cap B)$.
\end{proof}

\subsection{Boundary (边界)}\label{def:boundary}
\begin{definition}
  The \textbf{boundary} of $A$ is $\partial A = \bar{A} \setminus
  \operatorname{int}(A)$.
\end{definition}
Equivalently: $x \in \partial A \iff$ every neighborhood of $x$ meets
both $A$ and $X \setminus A$.

\begin{proof}
  ($\Rightarrow$) Let $x \in \partial A = \bar{A} \setminus
  \operatorname{int}(A)$.
  Since $x \in \bar{A}$, every neighborhood of $x$ meets $A$.
  Since $x \notin \operatorname{int}(A)$, no neighborhood of $x$ is
  contained in $A$,
  so every neighborhood meets $X \setminus A$.

  ($\Leftarrow$) If every neighborhood of $x$ meets $A$, then $x \in \bar{A}$.
  If every neighborhood meets $X \setminus A$, then no neighborhood
  is contained in $A$,
  so $x \notin \operatorname{int}(A)$. Hence $x \in \bar{A} \setminus
  \operatorname{int}(A) = \partial A$.
\end{proof}

\subsection{Neighborhood (邻域)}\label{def:neighborhood}
A \textbf{neighborhood} of $x \in X$ is a set $N$ containing an open
set $U$ with $x \in U \subseteq N$.
(Some authors require $N$ to be open itself; the definitions are
equivalent for most purposes.)

\subsection{Limit Point / Cluster Point / Accumulation Point
(极限点/聚点/附着点)}\label{def:limit-point}
\begin{definition}
  $x \in X$ is a \textbf{limit point} of $A$ if for every open set $U \ni x$:
  \[
    (U \cap A) \setminus \{x\} \ne \varnothing.
  \]
  I.e., every neighborhood of $x$ contains a point of $A$ different from $x$.
\end{definition}
\textbf{Examples:}
\begin{itemize}
  \item $A = \{1/n : n \in \N\}$: the only limit point is $0$ (and $0
    \notin A$).
  \item $A = \{a\}$ (singleton): no limit points.
\end{itemize}

\subsection{Isolated Point (孤立点)}\label{def:isolated}
$x \in A$ is an \textbf{isolated point} of $A$ if there exists an
open neighborhood $U$ of $x$ with $U \cap A = \{x\}$.
Equivalently: $x \in A$ is not a limit point of $A$.

\subsection{Derived Set (导集)}\label{def:derived}
The \textbf{derived set} of $A$, denoted $A'$ or $A^d$, is the set of
all limit points of $A$.
\textbf{Key relation:} $\bar{A} = A \cup A'$.

\subsection{Dense Set (稠密集)}\label{def:dense}
$A$ is \textbf{dense} in $X$ if $\bar{A} = X$. Equivalently: $A$
meets every nonempty open set.
Example: $\Q$ is dense in $\R$.

\subsection{Nowhere Dense (疏集/无处稠密)}\label{def:nowhere-dense}
$A$ is \textbf{nowhere dense} if $\operatorname{int}(\bar{A}) = \varnothing$.
Example: $\Z$ is nowhere dense in $\R$.

\subsection{Regularly Open (正则开集)}\label{def:regularly-open}
\begin{definition}
  $A$ is \textbf{regularly open} if $A = (\bar{A})^\circ$, i.e., $A$
  equals the interior of its closure.
\end{definition}
\textbf{Example:} $(0, 1)$ is regularly open; $(0, 1) \cup (1, 2)$ is
open but \emph{not} regularly open (its closure is $[0, 2]$, interior
of closure is $(0, 2) \ne (0, 1) \cup (1, 2)$).

\begin{proposition}
  For any $A \subseteq X$, the set $(\bar{A})^\circ$ is regularly open.
\end{proposition}

%=====================================================================
\section{连续映射 (Continuous Functions)}
%=====================================================================

\subsection{Continuous Function (连续函数/连续映射)}\label{def:continuous}
\begin{theorem}[Four equivalent definitions]
  Let $f : X \to Y$ where $X, Y$ are topological spaces. The
  following are equivalent:
  \begin{enumerate}
    \item $f^{-1}(V)$ is open in $X$ for every open $V \subseteq Y$
      (global continuity).
    \item For every $x \in X$ and every open $V \ni f(x)$, there
      exists open $U \ni x$ with $f(U) \subseteq V$ (continuity at each point).
    \item $f^{-1}(F)$ is closed in $X$ for every closed $F \subseteq Y$.
    \item $f(\bar{A}) \subseteq \overline{f(A)}$ for all $A \subseteq X$.
  \end{enumerate}
\end{theorem}

\begin{proof}[Proof of equivalence]
  (1)$\Rightarrow$(2): For $x \in X$ and open $V \ni f(x)$, set $U = f^{-1}(V)$.
  Then $U$ is open by (1), $x \in U$, and $f(U) \subseteq V$.

  (2)$\Rightarrow$(3): Let $F \subseteq Y$ be closed. To show
  $f^{-1}(F)$ is closed,
  take $x \in \overline{f^{-1}(F)}$. For every open $U \ni x$, $U
  \cap f^{-1}(F) \ne \varnothing$,
  so $\exists\, y \in U$ with $f(y) \in F$.
  By (2), $\exists$ open $U_0 \ni x$ with $f(U_0) \subseteq Y \setminus F$ open.
  But $U_0 \cap f^{-1}(F) \ne \varnothing$ gives a point mapped into
  $F$ and not into $F$, contradiction.
  Hence $x \in f^{-1}(F)$, so $f^{-1}(F)$ is closed.

  (3)$\Rightarrow$(4): Let $A \subseteq X$. Then $f(\bar{A})
  \subseteq \overline{f(A)}$ iff
  $\bar{A} \subseteq f^{-1}(\overline{f(A)})$. Since
  $\overline{f(A)}$ is closed, by (3),
  $f^{-1}(\overline{f(A)})$ is closed. Since $A \subseteq
  f^{-1}(\overline{f(A)})$, and $\bar{A}$
  is the smallest closed superset of $A$, we get $\bar{A} \subseteq
  f^{-1}(\overline{f(A)})$.

  (4)$\Rightarrow$(1): Let $V \subseteq Y$ be open. Set $A = X
  \setminus f^{-1}(V) = f^{-1}(Y \setminus V)$.
  Then $f(\bar{A}) \subseteq \overline{f(A)} \subseteq \overline{Y
  \setminus V} = Y \setminus V$ (closed).
  So $\bar{A} \subseteq f^{-1}(Y \setminus V) = A$, meaning $A$ is
  closed, hence $f^{-1}(V)$ is open.
\end{proof}

In metric spaces, definition (2) reduces to the
$\varepsilon$-$\delta$ formulation:
\[
  \forall \varepsilon > 0,\ \exists \delta > 0 :\ d(x, x_0) < \delta
  \implies d'(f(x), f(x_0)) < \varepsilon.
\]

\subsection{Homeomorphism (同胚)}\label{def:homeomorphism}
\begin{definition}
  $f : X \to Y$ is a \textbf{homeomorphism} if $f$ is a bijection,
  $f$ is continuous, and $f^{-1}$ is continuous.
\end{definition}
Homeomorphic spaces ($X \cong Y$) are topologically ``the same.''
\textbf{Counterexample:} $f : [0,1) \to S^1$, $f(t) = (\cos 2\pi t,
\sin 2\pi t)$ is a continuous bijection, but $f^{-1}$ is not continuous.

\subsubsection*{Technique 1: Construct an explicit homeomorphism}
To prove $X \cong Y$, build an explicit $f : X \to Y$ and verify:
\begin{enumerate}
  \item $f$ is a bijection (injective $+$ surjective, or exhibit
    $f^{-1}$ explicitly).
  \item $f$ is continuous.
  \item $f^{-1}$ is continuous.
\end{enumerate}
\textbf{Shortcut:} If $f$ is a bijection and $X$ is \emph{compact},
then ``$f$ continuous $\Rightarrow$ $f^{-1}$ continuous'' automatically.
Indeed, for $U$ open in $X$, $(f^{-1})^{-1}(U) = f(U)$; since $X
\setminus U$ is closed hence compact, $f(X \setminus U)$ is compact
hence closed, so $f(U) = Y \setminus f(X \setminus U)$ is open.

\textbf{Example:} $(0,1) \cong \R$ via $f(x) = \tan(\pi(x - 1/2))$.
Both $f$ and $f^{-1}(y) = \frac{1}{\pi}\arctan(y) + \frac{1}{2}$ are
continuous.

\subsubsection*{Technique 2: Show $X \not\cong Y$ via topological invariants}
Find a topological property that $X$ has but $Y$ does not (or vice
versa). Common invariants:
\begin{itemize}
  \item \textbf{Compactness:} $[0,1]$ is compact; $(0,1)$ and $(0,\infty)$
    are not. Hence $[0,1] \not\cong (0,1)$.
  \item \textbf{Connectedness:} $\R \setminus \{0\}$ is disconnected;
    $\R$ is connected. Hence $\R \not\cong \R \setminus \{0\}$.
  \item \textbf{Cardinality:} A finite space cannot be homeomorphic to an
    infinite one.
  \item \textbf{Hausdorff / separation axioms:} The trivial topology on
    $\{a,b\}$ is not Hausdorff; the discrete topology is. Not homeomorphic.
\end{itemize}

\subsubsection*{Technique 3: Point-removal argument}
If $X \cong Y$ via $f$, then for any $p \in X$:
\[
  X \setminus \{p\} \;\cong\; Y \setminus \{f(p)\}.
\]
So the topological properties of $X \setminus \{p\}$ must match those of
$Y \setminus \{f(p)\}$.

\textbf{Example:} $(0,1] \not\cong [0,1]$.
\begin{itemize}
  \item $(0,1] \setminus \{1\} = (0,1)$ is connected.
  \item $[0,1] \setminus \{p\}$ is disconnected for any $p \in (0,1)$.
  \item If a homeomorphism existed, $f(p)$ would need to be an endpoint,
    but then only $f(0)$ or $f(1)$ could map to endpoints, leaving no
    valid image for the other endpoint.
\end{itemize}

\textbf{Example:} $(0,\infty) \not\cong (0,1]$.
$(0,\infty) \setminus \{p\}$ is disconnected for every $p$, but
$(0,1] \setminus \{1\} = (0,1)$ is connected.

\textbf{Example:} $\R \not\cong \R^2$ (point-removal dimension
invariant): $\R \setminus \{p\}$ is disconnected; $\R^2 \setminus \{p\}$
is connected.

\subsubsection*{Technique 4: Open-mapping / closed-mapping criterion}
$f$ is a homeomorphism $\iff$ $f$ is a bijection, continuous, and
\textbf{open} (or equivalently \textbf{closed}).
\begin{itemize}
  \item $f$ is \textbf{open} if $f(U)$ is open in $Y$ for every open
    $U \subseteq X$.
  \item $f$ is \textbf{closed} if $f(F)$ is closed in $Y$ for every
    closed $F \subseteq X$.
\end{itemize}
\textbf{Use case:} When $f$ has a natural algebraic form (e.g.\ linear
maps), it may be easier to verify that $f$ is open than to construct
$f^{-1}$ explicitly.

\subsubsection*{Technique 5: Continuous bijection from compact to Hausdorff}
\begin{theorem}
  If $X$ is compact, $Y$ is Hausdorff, and $f : X \to Y$ is a
  continuous bijection, then $f$ is a homeomorphism.
\end{theorem}
\begin{proof}
  It suffices to show $f$ is closed. If $F \subseteq X$ is closed, then
  $F$ is compact (closed subset of compact). So $f(F)$ is compact
  (continuous image of compact). Since $Y$ is Hausdorff, $f(F)$ is
  closed.
\end{proof}
\textbf{Application:} Any continuous bijection from $[a,b]$ to a compact
Hausdorff space is automatically a homeomorphism.

\subsubsection*{Technique 6: Show a continuous bijection is NOT a homeomorphism}
To show $f$ is not a homeomorphism despite being a continuous bijection,
prove $f^{-1}$ is not continuous by finding an open set whose preimage
under $f^{-1}$ (i.e.\ image under $f$) is not open.

\textbf{Example:} $f : [0,1) \to S^1$, $f(t) = (\cos 2\pi t, \sin 2\pi t)$.
The set $[0, 1/2)$ is open in $[0,1)$, but $f([0,1/2))$ is not open in
$S^1$ (it contains the point $f(0) = (1,0)$ but no neighborhood of
$(1,0)$ in $S^1$).

\textbf{Alternative:} Use topological invariants. Here $S^1$ is compact
but $[0,1)$ is not.

\subsection{Topological Invariant (拓扑不变量)}\label{def:topo-invariant}
A property $P$ is a \textbf{topological invariant} if $X \cong Y$ and
$X$ has $P$ implies $Y$ has $P$.
Examples: connectedness, compactness, Hausdorff, cardinality.
Non-example: ``being a subset of $\R$'' is not a topological invariant.

\subsection{Embedding (嵌入)}\label{def:embedding}
$f : X \to Y$ is an \textbf{embedding} if $f : X \to f(X)$ is a
homeomorphism (where $f(X)$ has the subspace topology).

\subsection{Pasting / Gluing Lemma (粘接引理)}\label{def:pasting}
\begin{theorem}
  Let $X = A \cup B$ where $A, B$ are both closed (or both open). If
  $f|_A$ and $f|_B$ are continuous and agree on $A \cap B$, then $f$
  is continuous on $X$.
\end{theorem}
\begin{proof}
  (Closed case.) Check $f^{-1}(F)$ is closed for closed $F \subseteq Y$.
  $f^{-1}(F) = (f|_A)^{-1}(F) \cup (f|_B)^{-1}(F)$.
  Each $(f|_A)^{-1}(F)$ is closed in $A$ (continuity of $f|_A$),
  hence closed in $X$ (closed subset of closed set $A$).
  A finite union of closed sets is closed.

  (Open case.) For open $V \subseteq Y$: $f^{-1}(V) = (f|_A)^{-1}(V)
  \cup (f|_B)^{-1}(V)$.
  Each $(f|_A)^{-1}(V)$ is open in $A$, i.e.\ $A \cap W$ with $W$ open in $X$.
  Since $A$ is open, $A \cap W$ is open in $X$. Union of two open sets is open.
\end{proof}

\subsection{Restriction (限制)}\label{def:restriction}
$f|_A : A \to Y$ is the function $f|_A(x) = f(x)$ for $x \in A$. If
$f$ is continuous, so is $f|_A$.

\subsection{Composition (复合)}\label{def:composition}
If $f : X \to Y$ and $g : Y \to Z$ are continuous, then $g \circ f :
X \to Z$ is continuous.

\subsection{Coordinate Function (坐标函数)}\label{def:coordinate}
For $f : Y \to \prod_{i \in I} X_i$, the $i$-th coordinate function
is $f_i = \pi_i \circ f : Y \to X_i$.

%=====================================================================
\section{积拓扑与箱拓扑 (Product and Box Topology)}
%=====================================================================

\subsection{Motivation: Topology on a Product of Spaces}
Given topological spaces $\{X_i\}_{i \in I}$, the Cartesian product
$X = \prod_{i \in I} X_i$ is the set of all functions $x : I \to
\bigcup_{i} X_i$ with $x(i) \in X_i$ (written $(x_i)_{i \in I}$).

\textbf{Question:} What topology should $X$ carry?

There are two natural answers --- the \emph{product topology} and the
\emph{box topology}. For \textbf{finite} $I$ they coincide, but for
\textbf{infinite} $I$ they differ dramatically.

\subsection{Product Topology (积拓扑)}\label{def:product-topo}
\begin{definition}
  The \textbf{product topology} on $X = \prod_{i \in I} X_i$ has
  \textbf{subbasis}:
  \[
    \mathcal{S} = \{\pi_i^{-1}(U_i) : i \in I,\ U_i \text{ open in }
    X_i\},
  \]
  i.e.\ sets of the form $U_i \times \prod_{j \ne i} X_j$ (one
  coordinate restricted, all others free).

  The \textbf{basis} generated by $\mathcal{S}$ consists of finite
  intersections of subbasis elements:
  \[
    \mathcal{B} = \left\{\prod_{i \in I} U_i : U_i \text{ open in }
    X_i,\ U_i = X_i \text{ for all but finitely many } i\right\}.
  \]
\end{definition}

\textbf{Intuition:} A basic open set restricts only \emph{finitely
many} coordinates. The remaining coordinates are completely free (set to
$X_i$). Think of it as ``almost nothing is constrained.''

\textbf{Why ``for all but finitely many''?} Because the basis is formed
by \emph{finite intersections} of subbasis elements. Each subbasis
element restricts exactly one coordinate; intersecting finitely many of
them restricts finitely many coordinates.

\textbf{Consequence:} If $I$ is finite (say $|I| = n$), then ``all but
finitely many $=$ all,'' so the product topology and box topology are
the same.

\subsection{Box Topology (箱拓扑)}\label{def:box-topo}
\begin{definition}
  The \textbf{box topology} on $X = \prod_{i \in I} X_i$ has basis:
  \[
    \mathcal{B}_{\mathrm{box}} = \left\{\prod_{i \in I} U_i : U_i
    \text{ open in } X_i\right\}.
  \]
\end{definition}

\textbf{Intuition:} A basic open set restricts \emph{every} coordinate
independently. For each $i$, you choose an open $U_i \subseteq X_i$, and
the basic open set is the ``box'' $\prod_i U_i$.

\textbf{Comparison of bases (for infinite $I$):}
\[
  \mathcal{B}_{\mathrm{prod}} \subsetneq \mathcal{B}_{\mathrm{box}}.
\]
Every product-basic open set is also box-basic (set the ``free''
coordinates to $X_i$). But the box topology \emph{strictly} has more
basic open sets: a box with $U_i \ne X_i$ for infinitely many $i$ is
open in the box topology but \textbf{not} in the product topology.

\textbf{Therefore:} $\tau_{\mathrm{prod}} \subsetneq
\tau_{\mathrm{box}}$ for infinite $I$. The box topology is strictly
finer.

\subsection{Side-by-Side Comparison}

\begin{center}
  \begin{tabular}{@{}p{3.5cm}p{5.5cm}p{5.5cm}@{}}
    \hline
    & \textbf{Product Topology} & \textbf{Box Topology} \\
    \hline
    Basic open set & $\prod U_i$, $U_i = X_i$ for all but finitely many $i$ &
    $\prod U_i$, each $U_i$ open in $X_i$ \\[4pt]
    Subbasis & $\pi_i^{-1}(U_i)$ (one coordinate) &
    No simpler subbasis in general \\[4pt]
    Finer / coarser & Coarser & Strictly finer (for $|I| = \infty$) \\[4pt]
    Finite $I$ & $=$ box topology & $=$ product topology \\[4pt]
    Projections $\pi_i$ & Continuous \textbf{and open} & Continuous
    \textbf{and open} \\[4pt]
    Universal property & $f$ continuous $\iff$ each $f_i$ continuous &
    \textbf{Fails} \\[4pt]
    Continuous maps into & Easy to construct (coordinatewise) & Very
    restrictive \\[4pt]
    $\R^\omega$ connected? & Yes (path connected) & No (disconnected) \\
    $\R^\omega$ metrizable? & Yes & No \\
    $\R^\omega_{\mathrm{ez}}$ closure & All of $\R^\omega$ (dense) &
    $\R^\omega_{\mathrm{ez}}$ itself (closed) \\
    \hline
  \end{tabular}
\end{center}

\subsection{Projection Map (投影映射)}\label{def:projection}
$\pi_i : \prod_{j \in I} X_j \to X_i$, defined by $\pi_i((x_j)_{j \in
I}) = x_i$.

\textbf{Properties in both topologies:}
\begin{itemize}
  \item $\pi_i$ is continuous: $\pi_i^{-1}(U_i) = U_i \times \prod_{j
    \ne i} X_j$ is a subbasis element (product) / basic open set (box).
  \item $\pi_i$ is open: $\pi_i(\prod_j U_j) = U_i$ is open.
\end{itemize}

\textbf{Key point:} Projections are the same in both topologies. The
difference is in the \emph{other} direction --- maps \emph{into} the
product.

\subsection{Cylinder / Cylinder Set (柱集)}\label{def:cylinder}
A subbasis element of the product topology: $\pi_i^{-1}(U_i) = U_i
\times \prod_{j \ne i} X_j$.

\textbf{Visualization:} In $\R^3 = \R \times \R \times \R$, a cylinder
$(a,b) \times \R \times \R$ restricts only the first coordinate. A
basic open set in the product topology is a finite intersection of
cylinders, giving an open rectangle $(a_1,b_1) \times (a_2,b_2) \times
(a_3,b_3)$.

For infinite products, cylinders restrict only \emph{one} coordinate,
and basic open sets restrict \emph{finitely many}.

\subsection{Universal Property (泛性质)}\label{def:universal}
\begin{theorem}
  Let $X = \prod_{i \in I} X_i$ with the product topology. For any
  space $Y$:
  \[
    f : Y \to X \text{ is continuous} \iff \text{each } f_i = \pi_i
    \circ f : Y \to X_i \text{ is continuous.}
  \]
\end{theorem}
\begin{proof}
  ($\Rightarrow$) Each $\pi_i$ is continuous, and composition of
  continuous maps is continuous.

  ($\Leftarrow$) It suffices to check $f^{-1}(S)$ is open for every
  subbasis element $S = \pi_j^{-1}(U_j)$. Compute:
  \[
    f^{-1}(\pi_j^{-1}(U_j)) = (\pi_j \circ f)^{-1}(U_j) =
    f_j^{-1}(U_j),
  \]
  which is open since $f_j$ is continuous. Since $f^{-1}$ commutes with
  finite intersections and arbitrary unions, $f$ is continuous.
\end{proof}

\textbf{Why this fails for the box topology:} In the box topology, a
basic open set $\prod_i U_i$ with $U_i \ne X_i$ for infinitely many $i$
is NOT a finite intersection of subbasis elements. So checking on
subbasis elements is not enough. One can have all $f_i$ continuous but
$f$ not continuous in the box topology.

\textbf{Practical consequence:} In the product topology, you can define
maps into a product by specifying each coordinate independently. This
makes proofs much easier.

\subsection{Working with Product Topology: Practical Toolkit}

\textbf{How to show a set is open in the product topology:}
\begin{enumerate}
  \item Write it as a union of basic open sets $\prod U_i$ with $U_i =
    X_i$ for all but finitely many $i$.
  \item Or: show it is a union of finite intersections of cylinders
    $\pi_i^{-1}(U_i)$.
\end{enumerate}

\textbf{How to show a set is NOT open:}
\begin{enumerate}
  \item Assume it is open, so it contains a basic neighborhood of each
    of its points.
  \item Show that a basic neighborhood $\prod U_i$ (with $U_i = X_i$
    for all but finitely many $i$) must contain points outside the set.
  \item Exploit the \emph{free} coordinates: they can be chosen
    arbitrarily, so set them to values that leave the set.
\end{enumerate}

\textbf{Example:} $\R^\omega_{\mathrm{ez}}$ has empty interior in the
product topology. For any $x \in \R^\omega_{\mathrm{ez}}$ and any basic
neighborhood $\prod U_i$ (free for $i > N$), set $x_i = 1$ for some $i
> N$: this gives a point in the neighborhood that is not eventually
zero.

\subsection{Working with Box Topology: Practical Toolkit}

\textbf{How to show a set is open in the box topology:}
\begin{enumerate}
  \item For each point $x$ in the set, find open $U_i \subseteq X_i$
    with $x_i \in U_i$ for each $i$, such that $\prod_i U_i$ is
    contained in the set.
  \item This is easier than the product topology because you can control
    \emph{every} coordinate independently.
\end{enumerate}

\textbf{The ``box neighborhood trick'':} Given $x = (x_i)$, the set
$N_x = \prod_i (x_i - \varepsilon_i, x_i + \varepsilon_i)$ is a box
neighborhood. By choosing $\varepsilon_i$ appropriately, you can confine
all coordinates simultaneously.

\textbf{Example:} Bounded sequences form an open set in box topology.
If $|x_i| \le C$ for all $i$, then $\prod_i (x_i - 1, x_i + 1)$
contains only bounded sequences (each $y_i$ satisfies $|y_i| < C + 1$).

If $x$ is unbounded, then $\prod_i (x_i - 1, x_i + 1)$ contains only
unbounded sequences (each $y_i$ satisfies $|y_i| > |x_i| - 1$).

\subsection{Eventually-Zero Sequence (终零序列) --- A Key
Example}\label{def:eventually-zero}

Let $\R^\omega_{\mathrm{ez}} = \{(x_i) \in \R^\omega : x_i \ne 0
\text{ for only finitely many } i\}$.

\textbf{Interior in both topologies:} $\operatorname{int} = \varnothing$.

\emph{Product topology:} Any basic neighborhood of $x \in
\R^\omega_{\mathrm{ez}}$ has $U_i = \R$ for all $i > N$. For some $i_0
> N$ with $x_{i_0} = 0$, change coordinate $i_0$ to any nonzero value
--- still in the neighborhood, but not eventually zero.

\emph{Box topology:} For $x \in \R^\omega_{\mathrm{ez}}$, let $F =
\{i : x_i \ne 0\}$ (finite). Any box neighborhood $\prod U_i$ has
$0 \in U_i$ for $i \notin F$ (since $x_i = 0 \in U_i$). For $i
\notin F$, pick $t_i \in U_i \setminus \{0\}$ (possible since $U_i$ is
open in $\R$ and contains $0$). The point with coordinates $t_i$ for
$i \notin F$ and $x_i$ for $i \in F$ is in the box neighborhood but
has infinitely many nonzero coordinates.

\textbf{Closure in product topology:} $\overline{\R^\omega_{\mathrm{ez}}}
= \R^\omega$ (dense).

Given any $x \in \R^\omega$ and any basic neighborhood $\prod U_i$
(with $U_i = \R$ for $i > N$), define $y$ by $y_i = x_i$ for $i \le N$
and $y_i = 0$ for $i > N$. Then $y$ is eventually zero, and $y \in
\prod U_i$ (since $y_i = x_i \in U_i$ for $i \le N$, and $y_i = 0 \in
\R = U_i$ for $i > N$).

\textbf{Closure in box topology:}
$\overline{\R^\omega_{\mathrm{ez}}} =
\R^\omega_{\mathrm{ez}}$ (closed).

If $x \notin \R^\omega_{\mathrm{ez}}$, then $I_0 = \{i : x_i \ne 0\}$
is infinite. Define $\prod V_i$ by $V_i = \R \setminus \{0\}$ for $i
\in I_0$ and $V_i = \R$ for $i \notin I_0$. This is a box neighborhood
of $x$, and every point in it has nonzero $i$-th coordinate for all $i
\in I_0$, hence is not eventually zero. So $x \notin
\overline{\R^\omega_{\mathrm{ez}}}$.

\subsection{Bounded Sequence (有界序列)}\label{def:bounded-seq}
$U = \{x = (x_i) : \exists C_x > 0,\ |x_i| \le C_x\ \forall i\}$.

\textbf{In box topology,} $U$ and $\R^\omega \setminus U$ form a
nontrivial separation, proving $(\R^\omega, \mathcal{T}_b)$ is
disconnected.

\textbf{Proof that $U$ is box-open:} If $x \in U$ with bound $C$,
then $N_x = \prod_i (x_i - 1, x_i + 1) \subseteq U$ since $|y_i| \le
|x_i| + 1 \le C + 1$.

\textbf{Proof that $\R^\omega \setminus U$ is box-open:} If $x
\notin U$, then for every $n$ there exists $i_n$ with $|x_{i_n}| > n$.
Set $N_x = \prod_i (x_i - 1, x_i + 1)$. For any $y \in N_x$,
$|y_{i_n}| \ge |x_{i_n}| - 1 > n - 1$, which is unbounded as $n \to
\infty$. So $y \notin U$.

\subsection{Metrizability and Connectedness of $\R^\omega$}

\textbf{Product topology:} $(\R^\omega, \mathcal{T}_p)$ is
\textbf{metrizable} and \textbf{path connected}.

Explicit metric: $d_\omega(x,y) = \sup_{i \in \N}
\frac{d^*(x_i, y_i)}{i+1}$ where $d^*(u,v) = \min\{|u-v|, 1\}$.

Path connected: the straight-line homotopy $H(t) = ((1-t)x_i + ty_i)_i$
is continuous because each coordinate function $H_i(t) = (1-t)x_i +
ty_i$ is continuous, and the universal property guarantees continuity
in the product topology.

\textbf{Box topology:} $(\R^\omega, \mathcal{T}_b)$ is
\textbf{not metrizable} and \textbf{disconnected} (bounded/unbounded
separation above).

\begin{proposition}
  $(\R^\omega, \mathcal{T}_b)$ is not first-countable, hence not metrizable.
\end{proposition}
\begin{proof}
  Suppose $\mathbf{0} = (0,0,\ldots)$ has a countable neighborhood basis
  $\{U_n\}_{n \in \N}$. Each $U_n$ is a box neighborhood of $\mathbf{0}$,
  so $U_n$ contains a box of the form $\prod_{i=1}^\infty (-\varepsilon_{n,i},\,
  \varepsilon_{n,i})$ with $\varepsilon_{n,i} > 0$ for all $i$.
  Define $V = \prod_{i=1}^\infty (-\varepsilon_{i,i}/2,\; \varepsilon_{i,i}/2)$.
  Then $V$ is a box neighborhood of $\mathbf{0}$.
  But for each $n$, the $n$-th coordinate of $U_n$ extends to
  $\pm\varepsilon_{n,n}$,
  while $V$ restricts to $\pm\varepsilon_{n,n}/2$, so $U_n \not\subseteq V$.
  Hence no $U_n$ is contained in $V$, contradicting the assumption that
  $\{U_n\}$ is a neighborhood basis.
\end{proof}

\subsection{Torus (环面)}\label{def:torus}
$S^1 \times S^1$ (product of two circles). Visualized as the surface
of a donut. Constructed by gluing opposite edges of $[0,1] \times
[0,1]$.

%=====================================================================
\section{连通性 (Connectedness)}
%=====================================================================

\subsection{Connected Space (连通空间)}\label{def:connected}
\begin{definition}
  $(X, \mathcal{T})$ is \textbf{connected} if there do not exist
  nonempty disjoint open sets $U, V$ with $X = U \cup V$.
  Equivalently: the only clopen subsets of $X$ are $\varnothing$ and $X$.
\end{definition}

\subsection{Disconnected Space (不连通空间)}\label{def:disconnected}
$X$ is \textbf{disconnected} if there exists a \textbf{separation}:
nonempty open $U, V$ with $U \cap V = \varnothing$ and $U \cup V = X$.

\subsection{Totally Disconnected (完全不连通)}\label{def:totally-disc}
$X$ is \textbf{totally disconnected} if the only connected subsets
are singletons (and $\varnothing$).
Example: $\Q$ with the subspace topology from $\R$.

\subsection{Separation (分离)}\label{def:separation-conn}
A pair of nonempty, disjoint, open sets $U, V$ with $U \cup V = X$.

\subsection{Connected Component (连通分支)}\label{def:component}
Define $x \sim y$ if $x, y$ lie in a common connected subset. This is
an equivalence relation.
The equivalence class $C(x)$ is the \textbf{connected component} of
$x$: the maximal connected subset containing $x$.
\begin{itemize}
  \item Each component is connected.
  \item Different components are disjoint.
  \item Components are closed in $X$.
\end{itemize}

\begin{proof}[Proof that components are closed]
  Let $C$ be a connected component. Since $C$ is connected,
  $\overline{C}$ is connected
  (closure of connected). But $C$ is a \emph{maximal} connected set
  and $C \subseteq \overline{C}$,
  so $\overline{C} \subseteq C$, meaning $C = \overline{C}$,
  i.e.\ $C$ is closed.
\end{proof}

\subsection{Path Connected (道路连通)}\label{def:path-connected}
\begin{definition}
  $X$ is \textbf{path connected} if for every $a, b \in X$, there
  exists a continuous $f : [0, 1] \to X$ with $f(0) = a$ and $f(1) = b$.
\end{definition}

\begin{proposition}
  Path connected $\implies$ connected. The converse is \textbf{false}.
\end{proposition}
\begin{proof}
  Suppose $X$ is path connected but disconnected: $X = U \cup V$ is a
  separation.
  Pick $a \in U$, $b \in V$. By path connectedness, $\exists$
  continuous $\gamma : [0,1] \to X$
  with $\gamma(0) = a$, $\gamma(1) = b$. Then $\gamma^{-1}(U)$ and
  $\gamma^{-1}(V)$ are disjoint
  nonempty open subsets of $[0,1]$ with $\gamma^{-1}(U) \cup
  \gamma^{-1}(V) = [0,1]$.
  This is a separation of $[0,1]$, contradicting that $[0,1]$ is connected.
\end{proof}

\subsection{Path / Curve (道路/曲线)}\label{def:path}
A continuous map $f : [0, 1] \to X$ is called a \textbf{path} from
$f(0)$ to $f(1)$.

\subsection{Path Component (道路连通分支)}\label{def:path-component}
The equivalence class under ``$x \sim y$ if there is a path from $x$ to $y$.''
Each path component $\subseteq$ a connected component.

\subsection{Topologist's Sine Curve (拓扑学家的正弦曲线)}\label{def:sine-curve}
\[
  A = \{(x, \sin(1/x)) : x \in (0, 1]\} \cup \{(0, y) : y \in [-1, 1]\}.
\]
$A$ is \textbf{connected} but \textbf{not path connected}: no path
can reach $(0, y)$ from a point $(x, \sin(1/x))$ with $x > 0$.

\subsection{Intermediate Value Theorem --- IVT (介值定理)}\label{def:ivt}
\begin{theorem}
  If $f : [a, b] \to \R$ is continuous, then for every $c$ between
  $f(a)$ and $f(b)$, there exists $x \in [a, b]$ with $f(x) = c$.
\end{theorem}
\begin{proof}
  Suppose $c$ is not attained. Then $A = f^{-1}((-\infty, c))$ and $B
  = f^{-1}((c, +\infty))$ are nonempty open sets (preimages of open
  sets under continuous $f$) with $A \cup B = [a, b]$ and $A \cap B =
  \varnothing$. This contradicts the connectedness of $[a, b]$.
\end{proof}

\textbf{Topological form:} If $f : X \to \R$ is continuous and $X$ is
connected, then $f(X)$ is an interval (hence connected).

\subsection{Key Connectedness Theorems}\label{def:conn-thms}

\begin{theorem}[Continuous image of connected is connected]
  If $f : X \to Y$ is continuous and $X$ is connected, then $f(X)$ is connected.
\end{theorem}
\begin{proof}
  Contrapositive. Suppose $f(X)$ is disconnected: $f(X) = U \cup V$
  where $U, V$ are disjoint nonempty open sets in the subspace
  topology of $f(X)$.
  Then $f^{-1}(U)$ and $f^{-1}(V)$ are disjoint nonempty open sets in $X$
  (preimages of open sets under continuous $f$; nonempty because $f$
  is surjective onto $f(X)$),
  and $f^{-1}(U) \cup f^{-1}(V) = X$. This is a separation of $X$,
  contradiction.
\end{proof}

\begin{theorem}[Union theorem for connected sets]
  If $\{A_\alpha\}$ are connected subsets with
  $\bigcap_\alpha A_\alpha \ne \varnothing$, then $\bigcup_\alpha
  A_\alpha$ is connected.
\end{theorem}
\begin{proof}
  Pick $p \in \bigcap_\alpha A_\alpha$. Suppose $\bigcup A_\alpha = U \cup V$
  is a separation. Say $p \in U$. For each $\alpha$, $A_\alpha$ is connected,
  so $A_\alpha \subseteq U$ or $A_\alpha \subseteq V$ (otherwise
    $A_\alpha \cap U$,
  $A_\alpha \cap V$ would separate $A_\alpha$). Since $p \in A_\alpha \cap U$,
  we have $A_\alpha \subseteq U$ for all $\alpha$. Hence $\bigcup
  A_\alpha \subseteq U$,
  so $V = \varnothing$, contradiction.
\end{proof}

\begin{theorem}[Closure of connected is connected]
  If $A$ is connected, then $\bar{A}$ is connected.
\end{theorem}
\begin{proof}
  Suppose $\bar{A} = U \cup V$ is a separation with $U, V$ open in
  $\bar{A}$, disjoint, nonempty.
  Then $A = (A \cap U) \cup (A \cap V)$. Since $A$ is connected, one
  of these must be empty.
  Say $A \cap V = \varnothing$, so $A \subseteq U$.
  Now $V$ is a nonempty open subset of $\bar{A}$, so $V$ contains a
  point $x$ with
  every neighborhood of $x$ meeting $A \subseteq U$. But $V$ is a
  neighborhood of $x$ in $\bar{A}$,
  so $V \cap A \ne \varnothing$, contradicting $A \cap V = \varnothing$.
\end{proof}

\begin{theorem}[Connected subsets of $\R$ are exactly intervals]
  A subset $A \subseteq \R$ is connected if and only if $A$ is an interval
  (of any type: open, closed, half-open, unbounded, or a singleton).
\end{theorem}
\begin{proof}
  ($\Leftarrow$) Every interval $[a,b]$ is connected (standard proof
  via supremum; see T5 below).
  Half-open and open intervals are connected as subspaces of $[a,b]$
  or by direct argument.
  Unbounded intervals like $(-\infty, b)$ are unions of connected
  $[-n, b]$ with nonempty intersection.
  A singleton is trivially connected.

  ($\Rightarrow$) Suppose $A$ is not an interval. Then $\exists\, a,
  c \in A$ and $b \notin A$ with
  $a < b < c$. Set $U = A \cap (-\infty, b)$ and $V = A \cap (b, +\infty)$.
  Then $U, V$ are nonempty (contain $a$ and $c$ respectively), open in $A$,
  disjoint, and $U \cup V = A$. This is a separation.
\end{proof}

\begin{theorem}[$\R^n$ is connected]
  $\R^n$ is connected for all $n \ge 1$.
\end{theorem}
\begin{proof}
  Each line through the origin $\ell_v = \{tv : t \in \R\}$ is
  homeomorphic to $\R$, hence connected.
  All such lines pass through $0$. By the union theorem, $\R^n =
  \bigcup_v \ell_v$ is connected.
\end{proof}

\begin{theorem}[$\R^n \setminus \{0\}$ is connected for $n \ge 2$]
  For $n \ge 2$, $\R^n \setminus \{0\}$ is connected (in fact path connected).
\end{theorem}
\begin{proof}
  For any $x, y \in \R^n \setminus \{0\}$ with $n \ge 2$:
  if $x$ and $y$ are not antipodal (i.e., $y \ne -x/\|x\| \cdot \|y\|$),
  the straight line segment from $x$ to $y$ avoids $0$ (since the line through
  $x$ and $y$ would pass through $0$ only if they are antipodal).
  If $y = -cx$ for some $c > 0$, pick any $z$ not on the line through
  $x$ and $y$
  (possible since $n \ge 2$), and concatenate paths $x \to z \to y$,
  both avoiding $0$. Hence $\R^n \setminus \{0\}$ is path connected,
  hence connected.
\end{proof}

\begin{theorem}[Product of connected spaces is connected]
  If $\{X_i\}_{i \in I}$ are connected, then $\prod_{i \in I} X_i$
  (with the product topology) is connected.
\end{theorem}
\begin{proof}
  Fix $a = (a_i) \in \prod X_i$. For each finite $J \subseteq I$, define
  $A_J = \{x \in \prod X_i : x_i = a_i \text{ for all } i \notin J\}$.
  Then $A_J \cong \prod_{i \in J} X_i$ (a finite product of connected
    spaces is connected,
    by induction: $X_1 \times X_2$ is connected since $\{x_1\} \times
    X_2$ is connected for each $x_1$,
    and $(X_1 \times X_2) = \bigcup_{x_1} (\{x_1\} \times X_2)$ with
    all sets containing
  $(x_1, a_2)$ for fixed $a_2$).
  Each $A_J$ is connected and contains $a$. Let $A = \bigcup_{J
  \text{ finite}} A_J$.
  By the union theorem, $A$ is connected. Since $A$ is dense in $\prod X_i$
  (any basic open set restricts only finitely many coordinates, so
  contains a point of $A$),
  $\overline{A} = \prod X_i$ is connected.
\end{proof}

\subsection{Proof Techniques for Connectedness}

\subsubsection*{T1: Direct separation construction (to show disconnected)}
Construct two nonempty disjoint open sets $U, V$ with $U \cup V = X$.
\textbf{Equivalently,} find a nontrivial clopen subset ($\ne \varnothing, X$).

\textbf{Example:} $\Q$ is disconnected. For any irrational $a$, set
$U = \Q \cap (-\infty, a)$ and $V = \Q \cap (a, +\infty)$.

\subsubsection*{T2: Continuous function to a discrete space (to show
disconnected)}
Find a continuous surjection $f : X \to \{0, 1\}$ (discrete topology).
Then $f^{-1}(\{0\})$ and $f^{-1}(\{1\})$ form a separation.

\textbf{General principle:} Find continuous $f : X \to Y$ where $Y$ is
disconnected, then $X$ must be disconnected (contrapositive of
``continuous image of connected is connected'').

\textbf{Example:} $X = \{(x,0)\} \cup \{(x,1/x) : x > 0\} \subset \R^2$.
Define $f(x,y) = xy$. Then $f(X) = \{0, 1\}$, which is disconnected.
Hence $X$ is disconnected.

\textbf{Example:} Bounded and unbounded sequences in $(\R^\omega,
\mathcal{T}_b)$ are separated: the ``boundness'' function (though not a
real-valued function) can be replaced by observing that the two sets are
each open in box topology.

\subsubsection*{T3: Use known connected sets $+$ theorems (to show connected)}
Chain together known results rather than working from scratch:
\begin{enumerate}
  \item Start with a known connected set ($[0,1]$, $\R^n$, intervals).
  \item Apply: continuous image, closure, union with nonempty intersection,
    product.
\end{enumerate}

\textbf{Example:} $\R^n$ is connected: $\R^n = \bigcup_{v \in \R^n}
\ell_v$ where $\ell_v$ is the line through the origin in direction $v$.
Each line $\ell_v \cong \R$ is connected, and all lines pass through $0$,
so $\bigcup \ell_v$ is connected by the union theorem.

\textbf{Example:} The topologist's sine curve $\overline A$ is connected:
$A = \{(x, \sin(1/x)) : x \in (0,1]\}$ is connected (continuous image of
$(0,1]$), and $\overline{A}$ is connected (closure of connected).

\subsubsection*{T4: Clopen set argument (to show a property holds globally)}
To prove every point of a connected $X$ has property $P$:
\begin{enumerate}
  \item Define $S = \{x \in X : x \text{ has property } P\}$.
  \item Show $S$ is open (use local property or neighborhood argument).
  \item Show $S$ is closed (use limit points or complement is open).
  \item Since $X$ is connected and $S$ is clopen and nonempty, $S = X$.
\end{enumerate}

\textbf{Example:} Connected $+$ locally path connected $\Rightarrow$ path
connected. Set $S_x = \{y : \exists\text{ path from } x \text{ to } y\}$.
$S_x$ is open (path-connected neighborhood of $y$ is contained in $S_x$)
and closed (complement is open by same argument). So $S_x = X$.

\subsubsection*{T5: Show a subset of $\R$ is connected}
\textbf{Characterization:} $A \subseteq \R$ is connected $\iff$ $A$ is an
interval (any type: open, closed, half-open, unbounded, a single point).

\textbf{Proof that intervals are connected (standard proof):}
Suppose $[a,b] = U \cup V$ with $U, V$ open, disjoint, nonempty.
Pick $u \in U$, $v \in V$, say $u < v$. Set $c = \sup\{x \in U : x < v\}$.
Then $c \in [a,b]$. If $c \in U$: some neighborhood $(c - \varepsilon,
c + \varepsilon) \subseteq U$, but then $c + \varepsilon/2 \in U$ with
$c + \varepsilon/2 < v$, contradicting supremum. If $c \in V$: similar
contradiction.

\subsubsection*{T6: Show NOT connected --- by removing a point}
If removing a point disconnects $X$ but not $Y$, then $X \not\cong Y$
(as used in the homeomorphism section). More directly, if $X \setminus
\{p\}$ is disconnected for some $p$, this can sometimes prove $X$ itself
has a ``weakness'' at $p$.

\textbf{Example:} $\R \setminus \{0\} = (-\infty, 0) \cup (0, +\infty)$
is disconnected. Hence $\R \not\cong \R^2$ (since $\R^2 \setminus
\{(0,0)\}$ is connected).

\subsection{Proof Techniques for Path Connectedness}

\subsubsection*{P1: Construct explicit paths}
For each pair $a, b \in X$, write down a continuous $\gamma : [0,1] \to X$
with $\gamma(0) = a$, $\gamma(1) = b$.

\textbf{Common path constructions:}
\begin{itemize}
  \item \textbf{Straight line} (convex sets): $\gamma(t) = (1-t)a + tb$.
    Works in any convex subset of $\R^n$.
  \item \textbf{Piecewise linear:} go through intermediate waypoints.
  \item \textbf{Through a base point:} path $a \to x_0$, then $x_0 \to b$
    (concatenation). Only need to show every point has a path to $x_0$.
\end{itemize}

\textbf{Example:} $\R^n$ is path connected: straight line $\gamma(t) =
(1-t)\mathbf{a} + t\mathbf{b}$.

\textbf{Example:} $S^n$ ($n \ge 1$) is path connected: any two points
can be joined by an arc of a great circle, or go through a third point
not antipodal to either.

\subsubsection*{P2: Concatenation of paths}
If there is a path from $a$ to $c$ and a path from $c$ to $b$, then there
is a path from $a$ to $b$:
\[
  \gamma(t) =
  \begin{cases}
    \gamma_1(2t) & 0 \le t \le 1/2, \\
    \gamma_2(2t - 1) & 1/2 \le t \le 1.
  \end{cases}
\]
This is continuous by the pasting lemma ($[0,1/2]$ and $[1/2,1]$ are
closed in $[0,1]$, and $\gamma_1(1) = c = \gamma_2(0)$).

\textbf{Strategy:} Show every point has a path to a fixed base point
$x_0$; then any two points are connected via $x_0$.

\subsubsection*{P3: Continuous image preserves path connectedness}
\begin{proposition}
  If $X$ is path connected and $f : X \to Y$ is continuous and
  surjective, then $Y$ is path connected.
\end{proposition}
\begin{proof}
  Given $b_1, b_2 \in Y$, pick $a_i$ with $f(a_i) = b_i$. Let
  $\gamma$ be a path from $a_1$ to $a_2$ in $X$. Then $f \circ \gamma$
  is a path from $b_1$ to $b_2$ in $Y$.
\end{proof}

\textbf{Example:} $S^1$ is path connected (image of $[0,1]$ under
$t \mapsto (\cos 2\pi t, \sin 2\pi t)$). Any continuous surjective
image of an interval or $\R^n$ is path connected.

\subsubsection*{P4: Product of path connected spaces}
The product of path connected spaces is path connected (in product
topology). A path in $\prod X_i$ is given coordinatewise:
\[
  \gamma(t) = (\gamma_i(t))_{i \in I}
\]
where each $\gamma_i$ is a path in $X_i$. By the universal property,
$\gamma$ is continuous iff each coordinate is.

\textbf{Example:} $\R^\omega = \prod_{i \in \N} \R$ is path connected:
$H(t) = ((1-t)x_i + ty_i)_i$.

\subsubsection*{P5: Show NOT path connected --- exploit oscillation}
The standard counterexample is the \textbf{topologist's sine curve}:
\[
  A = \{(x, \sin(1/x)) : x \in (0,1]\}, \quad \overline{A} = A \cup
  (\{0\} \times [-1,1]).
\]
$\overline{A}$ is connected (closure of connected) but not path connected.

\textbf{Proof sketch (not path connected):} Suppose $\gamma : [0,1] \to
\overline{A}$ is a path from $(0,0)$ to $(1, \sin 1)$.
Write $\gamma(t) = (x(t), y(t))$. Then $x(t)$ is continuous, $x(0) = 0$,
$x(1) = 1$. Set $t_0 = \sup\{t : x(t) = 0\}$. For $t > t_0$,
$x(t) > 0$ so $y(t) = \sin(1/x(t))$. But $\sin(1/x)$ oscillates
infinitely as $x \to 0^+$, so $y$ cannot be continuous at $t_0$:
contradiction.

\textbf{Key insight:} When a function oscillates infinitely near a limit
point, no continuous path can ``trace'' the graph all the way to the
limit point.

\subsubsection*{P6: Connected $+$ locally path connected
$\Rightarrow$ path connected}
This is a powerful sufficient condition (see T4 above for the clopen proof).

\textbf{Examples where this applies:}
\begin{itemize}
  \item Open connected subsets of $\R^n$ are path connected (open balls
    are path connected, so $\R^n$ is locally path connected).
  \item Manifolds: connected manifolds are path connected.
\end{itemize}

\textbf{Warning:} ``Connected'' alone does NOT imply path connected
(topologist's sine curve). ``Locally path connected'' alone also does NOT
imply path connected (discrete spaces are locally path connected but may
have many path components). \emph{Both} hypotheses are needed.

%=====================================================================
\section{紧致性 (Compactness --- Introduction)}
%=====================================================================

\subsection{Compact Space (紧致空间/紧空间)}\label{def:compact}
\begin{definition}
  $(X, \mathcal{T})$ is \textbf{compact} if every open cover of $X$
  has a finite subcover.
\end{definition}

\subsection{Open Cover (开覆盖)}\label{def:open-cover}
A family $\{U_i\}_{i \in I}$ of open subsets is an \textbf{open
cover} of $X$ if $X = \bigcup_{i \in I} U_i$.

\subsection{Finite Subcover (有限子覆盖)}\label{def:finite-subcover}
A \textbf{finite subcover} is a finite subset $\{U_{i_1}, \ldots,
U_{i_k}\}$ that still covers $X$.

\subsection{Subcover (子覆盖)}\label{def:subcover}
A subfamily of a cover that still covers $X$.

\subsection{Heine--Borel Theorem (Heine--Borel 定理)}\label{def:heine-borel}
\begin{theorem}
  A subset $A \subseteq \R^n$ is compact (in the standard topology)
  if and only if $A$ is closed and bounded.
\end{theorem}
(See the Compactness Arguments section for the full proof.)

\textbf{Compact examples:} $[a, b]$; any finite space.
\textbf{Non-compact examples:} $\R$ (cover $\{(n, n+2)\}$ has no
finite subcover); $(0, 1)$; $\R^n$.

%=====================================================================
\section{贝尔纲定理 (Baire Category Theorem)}
%=====================================================================

\subsection{Nowhere Dense Set (疏朗集/无处稠密集)}\label{def:nowhere-dense-baire}
\begin{definition}
  A subset $A \subseteq X$ is \textbf{nowhere dense} if its closure has
  empty interior: $\operatorname{int}(\bar{A}) = \varnothing$.
\end{definition}

\textbf{Equivalent:} $A$ is nowhere dense iff $\bar{A}$ contains no
nonempty open set, i.e.\ $X \setminus \bar{A}$ is dense in $X$.

\textbf{Examples:}
\begin{itemize}
  \item $\Z \subseteq \R$: $\bar{\Z} = \Z$, $\operatorname{int}(\Z) =
    \varnothing$.
    Nowhere dense.
  \item $\Q \subseteq \R$: $\overline{\Q} = \R$,
    $\operatorname{int}(\R) = \R \ne \varnothing$.
    \textbf{Not} nowhere dense (even though $\Q$ has empty interior --- it's the
    \emph{closure} that matters).
  \item A single point $\{x\}$ in $\R$: nowhere dense.
  \item The Cantor set $C \subseteq [0,1]$: closed and has empty interior,
    hence nowhere dense.
  \item Any closed set with empty interior is nowhere dense.
\end{itemize}

\subsection{Meager / First Category Set (第一纲集)}\label{def:meager}
\begin{definition}
  $A \subseteq X$ is \textbf{meager} (or of \textbf{first category}) if
  $A$ is a countable union of nowhere dense sets:
  $A = \bigcup_{n=1}^\infty A_n$ where each $A_n$ is nowhere dense.
\end{definition}

\subsection{Non-meager / Second Category Set (第二纲集)}\label{def:non-meager}
\begin{definition}
  $A$ is of \textbf{second category} (in $X$) if it is \emph{not} meager
  (i.e.\ not a countable union of nowhere dense sets).
\end{definition}

\textbf{Warning:} ``First category'' $\neq$ ``small'' in a measure-theoretic
sense. For example, $\Q$ is meager in $\R$ but has full measure under
certain exotic measures. Conversely, the Cantor set is nowhere dense but
uncountable.

\subsection{Baire Space (Baire 空间)}\label{def:baire-space}
\begin{definition}
  A topological space $X$ is a \textbf{Baire space} if the intersection
  of any countable family of dense open sets is dense:
  \[
    \text{If } \{U_n\}_{n=1}^\infty \text{ are dense and open, then }
    \bigcap_{n=1}^\infty U_n \text{ is dense.}
  \]
\end{definition}

\textbf{Equivalent formulations} (for a topological space $X$):
\begin{enumerate}
  \item Every countable intersection of dense open sets is dense.
  \item Every countable union of closed nowhere dense sets has empty interior.
  \item Every nonempty open set is of second category (non-meager) in $X$.
\end{enumerate}

\begin{proof}[Equivalence ($1 \Leftrightarrow 2$)]
  Let $A_n$ be closed nowhere dense sets. Then $U_n = X \setminus A_n$ is
  open and dense (since $\operatorname{int}(A_n) = \varnothing$ means
  $X \setminus A_n$ is dense). Now:
  \[
    \operatorname{int}\!\left(\bigcup_{n=1}^\infty A_n\right) = \varnothing
    \;\iff\;
    X \setminus \bigcup_{n=1}^\infty A_n = \bigcap_{n=1}^\infty U_n
    \text{ is dense.}
  \]
\end{proof}

\subsection{Baire Category Theorem (贝尔纲定理)}\label{thm:baire}

\begin{theorem}[Baire Category Theorem]
  \begin{enumerate}
    \item Every \textbf{complete metric space} $(X, d)$ is a Baire space.
    \item Every \textbf{locally compact Hausdorff space} is a Baire space.
  \end{enumerate}
\end{theorem}

\begin{proof}[Proof (complete metric space case)]
  Let $\{U_n\}_{n=1}^\infty$ be dense open sets in $X$. We show
  $\bigcap_{n=1}^\infty U_n$ is dense by proving it meets every nonempty
  open ball $B_0 = B(x_0, r_0)$.

  Since $U_1$ is dense, $B_0 \cap U_1 \ne \varnothing$.
  Pick $x_1 \in B_0 \cap U_1$ and $r_1 < \min\{r_0/2, 1\}$ with
  $\bar{B}(x_1, r_1) \subseteq B_0 \cap U_1$.

  Inductively: pick $x_{n+1} \in B(x_n, r_n) \cap U_{n+1}$ (nonempty by density)
  and $r_{n+1} < \min\{r_n/2, 1/(n+1)\}$ with
  $\bar{B}(x_{n+1}, r_{n+1}) \subseteq B(x_n, r_n) \cap U_{n+1}$.

  The closed balls $\bar{B}(x_n, r_n)$ form a nested sequence with
  $\operatorname{diam}(\bar{B}(x_n, r_n)) \le 2r_n \to 0$.
  By \textbf{completeness}, $\bigcap_{n=1}^\infty \bar{B}(x_n, r_n) = \{x\}$
  for some $x \in X$.

  Then $x \in \bar{B}(x_n, r_n) \subseteq U_n$ for all $n$, so
  $x \in \bigcap_{n=1}^\infty U_n$.
  Also $x \in \bar{B}(x_1, r_1) \subseteq B_0$, so $x \in B_0$.
  Hence $B_0 \cap \bigcap_{n=1}^\infty U_n \ne \varnothing$.
\end{proof}

\begin{proof}[Proof (locally compact Hausdorff case)]
  Let $\{U_n\}$ be dense open sets and let $V$ be any nonempty open set.
  We show $V \cap \bigcap_n U_n \ne \varnothing$.

  Since $X$ is locally compact Hausdorff, there exists a nonempty open set
  $W_0$ with $\bar{W}_0$ compact and $\bar{W}_0 \subseteq V$.
  Since $U_1$ is dense, $W_0 \cap U_1 \ne \varnothing$; by local compactness
  again, pick open $W_1$ with $\bar{W}_1$ compact and
  $\bar{W}_1 \subseteq W_0 \cap U_1$.

  Inductively construct nested compact sets
  $\bar{W}_0 \supseteq \bar{W}_1 \supseteq \cdots$ with
  $\bar{W}_n \subseteq U_n$.
  By the \textbf{finite intersection property} of the compact set
  $\bar{W}_0$:
  $\bigcap_{n=0}^\infty \bar{W}_n \ne \varnothing$.
  Any point in this intersection lies in $V$ and in every $U_n$.
\end{proof}

\subsection{Applications of Baire Category Theorem
(贝尔纲定理的应用)}\label{def:baire-applications}

\textbf{Application 1: $\R$ is uncountable.}

$\R$ is a complete metric space, hence Baire. If $\R$ were countable,
say $\R = \{x_1, x_2, \ldots\}$, then $\R = \bigcup_{n=1}^\infty \{x_n\}$.
Each $\{x_n\}$ is nowhere dense (closed with empty interior).
But $\R$ itself has nonempty interior, contradicting Baire.
So $\R$ is uncountable.

\textbf{Application 2: $\Q$ is not a $G_\delta$ set in $\R$.}

A $G_\delta$ set is a countable intersection of open sets.
If $\Q = \bigcap_{n=1}^\infty U_n$ with each $U_n$ open,
then each $U_n$ is dense (since $\Q$ is dense), so by Baire,
$\bigcap_{n=1}^\infty U_n$ would be dense and hence uncountable.
But $\Q$ is countable --- contradiction.

\textbf{Application 3: Existence of continuous nowhere-differentiable
functions.}

The set of continuous functions on $[0,1]$ that are \emph{nowhere
differentiable}
is a dense $G_\delta$ in $C([0,1])$ (with the sup-norm metric), hence
\textbf{generic} (second category). In particular, ``most'' continuous
functions are nowhere differentiable --- they are the rule, not the exception.

\textbf{Application 4: $\R$ cannot be written as a countable union of
closed sets with empty interior.}

This is immediate from Baire: $\R$ has nonempty interior.

\subsection{Quick Reference: Baire Space Summary}\label{def:baire-summary}

\begin{center}
  \begin{tabular}{ll}
    \textbf{Space} & \textbf{Baire?} \\
    \hline
    $\R$, $\R^n$ (complete metric) & Yes \\
    Any complete metric space & Yes \\
    Any locally compact Hausdorff space & Yes \\
    $(0, 1) \subset \R$ (incomplete but locally compact Hausdorff) & Yes \\
    $\Q$ (with subspace topology) & \textbf{No} \\
    Any discrete space (locally compact Hausdorff) & Yes \\
  \end{tabular}
\end{center}

\textbf{Key ``mnemonic'':} Complete metric spaces and locally compact
Hausdorff spaces ``have enough room'' that countably many thin (nowhere dense)
sets cannot cover the whole space. You can always find a point that avoids
all of them.

%=====================================================================
\section{Key Results from Homework (HW1--3)}
%=====================================================================

\subsection{Characteristic Function Bijection (HW1 Q1)}
\begin{theorem}
  For any nonempty finite set $A$, there is a bijection between
  $\mathcal{P}(A)$ and $2^A = \{f : A \to \{0,1\}\}$.
  Hence $|\mathcal{P}(A)| = 2^{|A|}$.
\end{theorem}
\begin{proof}
  Define $\Phi : \mathcal{P}(A) \to 2^A$ by
  $\Phi(B) = \chi_B$ (the characteristic / indicator function of $B$):
  \[
    \chi_B(a) =
    \begin{cases} 1 & a \in B, \\ 0 & a \notin B.
    \end{cases}
  \]
  \textbf{Injective:} $\chi_{B_1} = \chi_{B_2} \Rightarrow B_1 = B_2$
  (check elementwise). \textbf{Surjective:} given $f \in 2^A$, set
  $B_f = f^{-1}(\{1\})$; then $\Phi(B_f) = f$.
\end{proof}
\textbf{Technique:} The characteristic function $\chi_A$ converts set
operations into pointwise arithmetic: $\chi_{A \cap B} = \chi_A \cdot
\chi_B$, $\chi_{A \cup B} = \max\{\chi_A, \chi_B\}$,
$\chi_{X \setminus A} = 1 - \chi_A$.

\subsection{Symmetric Group on a Finite Set (HW1 Q2)}
\begin{proposition}
  Let $F(A)$ be the set of all bijections $f : A \to A$. Then
  $(F(A), \circ)$ is a group. It is abelian if and only if $|A| \le 2$.
\end{proposition}
\textbf{Verification checklist} (group axioms):
\begin{enumerate}
  \item \textbf{Closure:} composition of bijections is a bijection.
  \item \textbf{Associativity:} $(f \circ g) \circ h = f \circ (g \circ
    h)$ for functions.
  \item \textbf{Identity:} $\mathrm{id}_A \in F(A)$.
  \item \textbf{Inverses:} every bijection $f$ has $f^{-1}$ with
    $f \circ f^{-1} = f^{-1} \circ f = \mathrm{id}$.
\end{enumerate}
\textbf{Non-abelian for $|A| \ge 3$:} Take $A = \{1,2,3\}$, $\sigma =
(1\;2)$, $\tau = (2\;3)$; then $\sigma \circ \tau \ne \tau \circ
\sigma$.

\subsection{Induction on Primes (HW1 Q3)}
\begin{theorem}
  If $p_n$ is the $n$-th prime, then $p_n < 2^{2^n}$ for all $n \ge 1$.
\end{theorem}
\begin{proof}
  Base: $p_1 = 2 < 2^2 = 4$.
  Inductive step: by Euclid's argument,
  \[
    p_{n+1} \le p_1 p_2 \cdots p_n + 1
    < 2^{2^1} \cdot 2^{2^2} \cdots 2^{2^n} + 1
    = 2^{2+4+\cdots+2^n} + 1 = 2^{2^{n+1}-2} + 1 < 2^{2^{n+1}}.
  \]
\end{proof}
\textbf{Technique:} When the inductive hypothesis gives a bound on a
product, use $\sum_{k=1}^n 2^k = 2^{n+1} - 2$ (geometric series).

\subsection{Equivalence Relations: Intersection and Generation (HW1 Q4)}
\begin{definition}
  A relation $R$ on $X$ is an \textbf{equivalence relation} if it is
  reflexive ($x R x$), symmetric ($x R y \Rightarrow y R x$), and
  transitive ($x R y, y R z \Rightarrow x R z$).
\end{definition}

\begin{proposition}
  The intersection $\bigcap_{i \in I} R_i$ of any family of equivalence
  relations on $X$ is an equivalence relation.
\end{proposition}
\begin{proof}
  Each axiom is checked pointwise. E.g., for transitivity: if
  $(x,y), (y,z) \in \bigcap R_i$, then $(x,y), (y,z) \in R_i$ for
  each $i$, so $(x,z) \in R_i$ for each $i$ (transitivity of $R_i$),
  hence $(x,z) \in \bigcap R_i$.
\end{proof}

\textbf{Generated equivalence relation:} The smallest equivalence
relation containing a given relation $S$ is the intersection of all
equivalence relations containing $S$. For the relation $S = \{(x,y) :
y = x+1, 0 < x < 3\}$ on $\R$, the generated equivalence relation is
$y - x \in \Z$ (since $S$ connects $x$ to $x+1$, transitivity gives all
integer differences).

\subsection{Countability of Algebraic Numbers (HW2 Q1)}
\begin{theorem}
  The set $A$ of real algebraic numbers is countable. Hence the set of
  real transcendental numbers $\R \setminus A$ is uncountable.
\end{theorem}
\begin{proof}
  A rational polynomial of degree $n$ is determined by $n+1$ rational
  coefficients $(a_0, \ldots, a_{n-1}, 1)$, so the set of such
  polynomials is countable ($\Q^n$ is countable).
  The set of \emph{all} rational polynomials is
  $\bigcup_{n=1}^\infty P_n$, a countable union of countable sets,
  hence countable.
  Each polynomial has finitely many roots, so $A$ is a countable union
  of finite sets, hence countable.
  Since $\R$ is uncountable, $\R \setminus A$ is uncountable.
\end{proof}
\textbf{Technique:} Countable union of countable (or finite) sets is
countable. Use this to reduce ``uncountable'' objects to countable
pieces.

\subsection{Eventually-Zero Binary Sequences (HW2 Q2)}
\begin{proposition}
  $X_0^\N = \{f : \N \to \{0,1\} : |\mathrm{supp}(f)| < \infty\}$
  is countable.
\end{proposition}
\begin{proof}
  Partition by support size: $X_0^\N = \bigcup_{n=0}^\infty S_n$ where
  $S_n = \{f : |\mathrm{supp}(f)| = n\}$.
  Each $S_n$ is countable (inject into $\N^n$ by listing support
  positions). A countable union of countable sets is countable.
\end{proof}
\textbf{Contrast with HW5:} When the support can be infinite, i.e.\
$\{0,1\}^\N$ (all binary sequences), the set is uncountable by Cantor
diagonalization.

\subsection{Pseudo-metric Quotient (HW2 Q3)}
\begin{definition}
  A \textbf{pseudo-metric} on $S$ is $d : S \times S \to \R_{\ge 0}$
  satisfying $d(x,x)=0$, symmetry, and the triangle inequality (but
  possibly $d(x,y)=0$ for $x \ne y$).
\end{definition}

\begin{proposition}
  Define $x \sim y$ iff $d(x,y) = 0$. Then:
  \begin{enumerate}
    \item $\sim$ is an equivalence relation.
    \item $d^*([x],[y]) = d(x,y)$ is a well-defined metric on $S^* = S/{\sim}$.
  \end{enumerate}
\end{proposition}
\textbf{Key technique --- well-definedness:} Must show $d^*$ does not
depend on the choice of representatives: if $x \sim x'$ and $y \sim y'$,
then $d(x,y) = d(x',y')$. Proof: $d(x,y) \le d(x,x') + d(x',y') +
d(y',y) = d(x',y')$, and symmetrically.

\subsection{Verifying the Triangle Inequality (HW2 Q4)}
\textbf{Problem:} For which $n \in \N_{\ge 1}$ is $d_n(x,y) =
(|x_1 - y_1| + |x_2 - y_2|)^n$ a metric on $\R^2$?

\textbf{Answer:} Only $n = 1$.
\begin{itemize}
  \item $n = 1$: $d_1$ is the $\ell^1$ metric; triangle inequality is standard.
  \item $n \ge 2$: Triangle inequality fails. Take $x = (0,0)$, $y =
    (1,0)$, $z = (1,1)$. Then $d_n(x,z) = 2^n$, but $d_n(x,y) + d_n(y,z)
    = 1 + 1 = 2 < 2^n$ for $n \ge 2$.
\end{itemize}
\textbf{Technique:} To disprove a metric axiom, find specific
counterexample points. The triangle inequality is usually the first to
fail.

\subsection{No Metric on Power Set Extending Point Distances (HW2 Q5)}
\begin{proposition}
  In general, there is no metric on $\mathcal{P}(S)$ such that
  $d'(\{x\},\{y\}) = d(x,y)$ for all $x, y \in S$.
\end{proposition}
\textbf{Counterexample:} Take $S = \{a, b, c\}$ with $d(a,b) = d(b,c)
= d(a,c) = 1$ (discrete metric). Set $d'(\{a\},\{b,c\}) = t$.
By triangle inequality: $t = d'(\{a\},\{b,c\}) \le d'(\{a\},\{b\}) +
d'(\{b\},\{b,c\})$, but this forces constraints that cannot be
simultaneously satisfied.

\textbf{Technique:} ``Can one always\ldots'' questions are often
resolved by a small counterexample ($|S| = 3$ suffices here).

\subsection{Discrete Metric Induces Discrete Topology (HW3 Q1)}
\begin{proposition}
  For any set $S$ with the discrete metric $d(x,y) = \mathbf{1}_{x \ne
  y}$, the induced topology is $\mathcal{P}(S)$.
\end{proposition}
\begin{proof}
  For each $x \in S$, the open ball $B(x, 1/2) = \{x\}$ is open.
  Any $A \subseteq S$ is a union of singletons: $A = \bigcup_{x \in A}
  \{x\}$, hence open.
\end{proof}

\textbf{Conversely:} The trivial topology $\{\varnothing, S\}$ is
metrizable only if $|S| \le 1$. If $|S| \ge 2$ and $d$ is a metric,
pick distinct $x, y$; then $B(x, d(x,y)/2)$ is a nonempty open set
$\ne S$, contradicting the trivial topology.

\subsection{Rational Intervals as a Countable Basis (HW3 Q2)}
\begin{proposition}
  $\mathcal{B} = \{(a,b) : a < b,\ a,b \in \Q\}$ is a basis for the
  standard topology on $\R$, and $|\mathcal{B}| = \aleph_0$.
\end{proposition}
\textbf{Basis verification:}
\begin{enumerate}
  \item \textbf{Covering:} $\bigcup_{a<b,\, a,b \in \Q} (a,b) = \R$
    (for any $x \in \R$, pick $a < x < b$ with $a, b \in \Q$ by
    density of $\Q$).
  \item \textbf{Intersection:} if $x \in (a_1,b_1) \cap (a_2,b_2)$,
    pick $c_1, c_2 \in \Q$ with $\max(a_1,a_2) < c_1 < x < c_2 <
    \min(b_1,b_2)$; then $x \in (c_1,c_2) \subseteq (a_1,b_1) \cap
    (a_2,b_2)$.
\end{enumerate}
\textbf{Countability:} injection $\mathcal{B} \hookrightarrow \Q \times
\Q$ via $(a,b) \mapsto (a,b)$. Since $\Q \times \Q$ is countable, so is
$\mathcal{B}$.

\subsection{Intersection of Topologies (HW3 Q4)}
\begin{proposition}
  Let $\{\mathcal{T}_i\}_{i \in I}$ be a family of topologies on $S$.
  Then $\bigcap_{i \in I} \mathcal{T}_i$ is a topology on $S$.
\end{proposition}
\begin{proof}
  \begin{enumerate}
    \item $\varnothing, S \in \mathcal{T}_i$ for all $i$, hence in the
      intersection.
    \item If $\{U_\alpha\} \subseteq \bigcap \mathcal{T}_i$, then each
      $U_\alpha \in \mathcal{T}_i$ for all $i$, so $\bigcup_\alpha
      U_\alpha \in \mathcal{T}_i$ for all $i$, hence in the
      intersection.
    \item Finite intersections: analogous.
  \end{enumerate}
\end{proof}

\textbf{Warning:} $\mathcal{T}_1 \cup \mathcal{T}_2$ need \emph{not}
be a topology. \textbf{Counterexample:} $S = \{a,b,c\}$,
$\mathcal{T}_1 = \{\varnothing, \{a\}, S\}$,
$\mathcal{T}_2 = \{\varnothing, \{b\}, S\}$.
Then $\{a\}, \{b\} \in \mathcal{T}_1 \cup \mathcal{T}_2$ but $\{a,b\}
\notin \mathcal{T}_1 \cup \mathcal{T}_2$.

\textbf{Generated topology:} The smallest topology containing
$\mathcal{T}_1 \cup \mathcal{T}_2$ is
$\bigcap\{\mathcal{T} : \mathcal{T}_1 \cup \mathcal{T}_2 \subseteq
\mathcal{T}\}$. Equivalently, take $\mathcal{T}_1 \cup \mathcal{T}_2$
as a \textbf{subbasis} and generate.

\subsection{Radially Open Topology (HW3 Q5)}
\begin{definition}
  $S \subseteq \R^2$ is \textbf{radially open} if for each $p \in S$
  and each direction (unit vector $v$), $S$ contains an open line
  segment from $p$ in direction $v$.
\end{definition}

\begin{proposition}
  $\mathcal{T}_{\mathrm{ro}}$ (all radially open sets) is a topology on
  $\R^2$, and $\mathcal{T}_{\mathrm{std}} \subsetneq
  \mathcal{T}_{\mathrm{ro}}$.
\end{proposition}
\textbf{Topology verification:} $\varnothing$ and $\R^2$ are trivially
radially open. Unions preserve radial openness (if each $S_\alpha$
contains a segment, so does the union). Finite intersections: at $p \in
\bigcap_{k=1}^n S_k$, take the minimum segment length over all $S_k$.

\textbf{Strictly finer:} Every open ball $B(p,r)$ is radially open
(ball contains segment in every direction). But a ``cross'' shape
$C = \{(x,0) : |x|<1\} \cup \{(0,y) : |y|<1\}$ is radially open
(segments exist along axes) but contains no open ball around the origin,
so $C \in \mathcal{T}_{\mathrm{ro}} \setminus \mathcal{T}_{\mathrm{std}}$.

%=====================================================================
\section{Key Results from Homework (HW4--6)}
%=====================================================================

\subsection{Closure Operator Axioms $\Rightarrow$ Topology (HW4 Q2)}
\begin{theorem}[Kuratowski closure axioms]
  Let $f : \mathcal{P}(X) \to \mathcal{P}(X)$ satisfy:
  \begin{enumerate}
    \item $f(\varnothing) = \varnothing$,
    \item $A \subseteq f(A)$ (expanding),
    \item $f(f(A)) = f(A)$ (idempotent),
    \item $f(A \cup B) = f(A) \cup f(B)$ (finite additivity).
  \end{enumerate}
  Then $\mathcal{C}_f = \{C \subseteq X : f(C) = C\}$ is the
  collection of closed sets for a topology $\mathcal{T}_f = \{X
  \setminus C : C \in \mathcal{C}_f\}$, and
  $\overline{A}^{\mathcal{T}_f} = f(A)$.
\end{theorem}
\textbf{Key technique --- monotonicity from additivity:} $A \subseteq
B \Rightarrow B = A \cup B \Rightarrow f(B) = f(A) \cup f(B)
\Rightarrow f(A) \subseteq f(B)$. This is used repeatedly.

\begin{proof}[Proof that $\mathcal{C}_f$ defines a topology]
  \textbf{Step 1: Monotonicity.} $A \subseteq B \Rightarrow f(A) \subseteq f(B)$
  (from $f(B) = f(A) \cup f(B \setminus A) \supseteq f(A)$).

  \textbf{Step 2: $\mathcal{C}_f$ contains $\varnothing$ and $X$.}
  $f(\varnothing) = \varnothing$ by axiom (1), so $\varnothing \in
  \mathcal{C}_f$.
  $X \subseteq f(X)$ by (2) and $f(X) \subseteq X$ trivially (since
  $f$ maps into $\mathcal{P}(X)$),
  so $f(X) = X$ and $X \in \mathcal{C}_f$.

  \textbf{Step 3: Finite unions.} If $C_1, C_2 \in \mathcal{C}_f$, then
  $f(C_1 \cup C_2) = f(C_1) \cup f(C_2) = C_1 \cup C_2$ by axiom (4),
  so $C_1 \cup C_2 \in \mathcal{C}_f$.

  \textbf{Step 4: Arbitrary intersections.} Let $\{C_\alpha\}
  \subseteq \mathcal{C}_f$,
  $C = \bigcap C_\alpha$. Then $C \subseteq C_\alpha$ implies $f(C)
  \subseteq f(C_\alpha) = C_\alpha$
  by monotonicity, so $f(C) \subseteq \bigcap C_\alpha = C$. Combined
  with $C \subseteq f(C)$,
  $f(C) = C$, hence $C \in \mathcal{C}_f$.

  So $\mathcal{C}_f$ is a collection of closed sets. Hence
  $\mathcal{T}_f = \{X \setminus C : C \in \mathcal{C}_f\}$ is a topology.

  \textbf{Step 5: $f(A) = \overline{A}^{\mathcal{T}_f}$.}
  Since $A \subseteq f(A) = f(f(A))$, $f(A) \in \mathcal{C}_f$, so
  $f(A)$ is closed.
  If $A \subseteq C \in \mathcal{C}_f$, then $f(A) \subseteq f(C) =
  C$ by monotonicity.
  So $f(A)$ is the smallest closed superset of $A$, i.e.\ $f(A) = \overline{A}$.
\end{proof}

\subsection{Closure $=$ Set $\cup$ Limit Points (HW4 Q5)}
\begin{theorem}
  $\overline{A} = A \cup LP(A)$.
\end{theorem}
\begin{proof}
  ($\subseteq$) Take $x \in \overline{A}$. If $x \in A$, done. If $x
  \notin A$: every neighborhood $U$ of $x$ meets $A$ (since $x \in
  \overline{A}$), and $x \notin A$, so $U \cap (A \setminus \{x\})
  \ne \varnothing$. Hence $x \in LP(A)$.

  ($\supseteq$) $A \subseteq \overline{A}$ trivially. If $x \in
  LP(A)$: every neighborhood of $x$ meets $A \setminus \{x\}
  \subseteq A$, so every neighborhood meets $A$, hence $x \in \overline{A}$.
\end{proof}

\subsection{Computing Interior, Closure, Limit Points (HW4 Q1)}
\textbf{Example (topologist's sine curve):}
$A = \{(x, \sin(1/x)) : x \in (0, 1]\} \subset \R^2$.
\begin{itemize}
  \item $A^\circ = \varnothing$ (graph of a function over $(0,1]$ has
    empty interior in $\R^2$).
  \item $\overline{A} = A \cup (\{0\} \times [-1, 1])$.
  \item $LP(A) = A \cup (\{0\} \times [-1, 1])$ (every point of $A$
    and the vertical segment are limit points).
  \item $\overline{A} \cap (\R^2 \setminus A) = \{0\} \times [-1, 1]$
    (the ``accumulated boundary'').
\end{itemize}
\textbf{Proof technique --- sequence construction:} For each $y \in
[-1, 1]$, choose $t_n \to \infty$ with $\sin t_n \to y$, set $x_n =
1/t_n$. Then $(x_n, \sin(1/x_n)) = (x_n, \sin t_n) \to (0, y)$.

\subsection{Compactness as a Topological Invariant (HW4 Q3)}
$f : [0,1) \to S^1$, $f(t) = (\cos 2\pi t, \sin 2\pi t)$ is a
continuous bijection but NOT a homeomorphism.
\textbf{Reason:} $S^1$ is compact (closed and bounded in $\R^2$);
$[0,1)$ is not compact. Compactness is a topological invariant.
\textbf{Consequence:} No homeomorphism exists between $[0,1)$ and $S^1$ at all.

\subsection{Regularly Open Sets (HW4 Q4)}
\begin{definition}
  $A$ is \textbf{regularly open} if $A = (\overline{A})^\circ$.
\end{definition}
\textbf{Example of open but not regularly open:} $U = (0,1) \cup (1,2)$ in $\R$:
$\overline{U} = [0,2]$, $(\overline{U})^\circ = (0,2) \ne U$.

\textbf{Characterization:} $U$ is regularly open $\iff$ $U =
\operatorname{int}(F)$ for some closed $F$.

\begin{proposition}
  For any $A \subseteq X$, $(\overline{A})^\circ$ is regularly open.
\end{proposition}
\begin{proof}
  Set $B = (\overline{A})^\circ$. Then $B$ open $\Rightarrow B
  \subseteq (\overline{B})^\circ$.
  Also $B \subseteq \overline{A} \Rightarrow \overline{B} \subseteq
  \overline{\overline{A}} = \overline{A}$, so $(\overline{B})^\circ
  \subseteq (\overline{A})^\circ = B$.
  Hence $(\overline{B})^\circ = B$.
\end{proof}

\subsection{Bounded Metric Trick (HW5 Q1)}
\begin{proposition}
  $d^*(x,y) = \min\{d(x,y), 1\}$ is a metric inducing the same topology as $d$.
\end{proposition}
\textbf{Triangle inequality --- case analysis:} Need $\min\{a+b, 1\}
\le \min\{a,1\} + \min\{b,1\}$.
If $a \ge 1$ or $b \ge 1$: RHS $\ge 1 \ge$ LHS. If $a < 1, b < 1$:
RHS $= a+b \ge$ LHS.

\textbf{Same topology --- ball comparison:} For $0 < r \le 1$:
$B_{d^*}(x,r) = B_d(x,r)$. For general balls, use $r' = \min\{r, 1/2\}$.

\textbf{Warning:} No uniform $A > 0$ with $A \cdot d \le d^*$ in
general (take $X = \R$).

\subsection{Free Coordinates in Product/Box Topology (HW5 Q2)}
\textbf{Key technique:} In the product topology, ``free''
(unrestricted) coordinates are a powerful tool.

\textbf{Eventually-zero sequences} $\R^\omega_{\mathrm{ez}}$ in $\R^\omega$:
\begin{itemize}
  \item Interior is $\varnothing$ in \emph{both} topologies: set free
    coordinates to $1$ (product) or replace zeros with $t_i \ne 0$ (box).
  \item Closure in $\mathcal{T}_p$: all of $\R^\omega$.
    \textbf{Reason:} given $x \in \R^\omega$ and a basic
    neighborhood, only finitely many coordinates are restricted; set
    the rest to $0$.
  \item Closure in $\mathcal{T}_b$: $\R^\omega_{\mathrm{ez}}$ itself
    (it is closed). \textbf{Reason:} if $x \notin
    \R^\omega_{\mathrm{ez}}$, then $I = \{i : x_i \ne 0\}$ is
    infinite; box-neighborhood $\prod_i V_i$ with $0 \notin V_i$ for
    $i \in I$ stays in the complement.
\end{itemize}

\subsection{Metrizability of $\R^\omega$ (HW5 Q3)}
\begin{theorem}
  $(\R^\omega, \mathcal{T}_p)$ is metrizable and connected.
\end{theorem}
\textbf{Explicit metric:} $d_\omega(x,y) = \sup_{i \in \N}
\frac{d^*(x_i, y_i)}{i+1}$ where $d^*(u,v) = \min\{|u-v|, 1\}$.

\textbf{Topology comparison --- two directions:}
\begin{itemize}
  \item $d_\omega$-ball $\subseteq$ $\mathcal{T}_p$-open: Choose $N$
    with $1/(N+1) < \varepsilon/2$; restrict first $N$ coordinates.
    Tail automatically satisfies $d^*/(i+1) < \varepsilon/2$ for $i > N$.
  \item $\mathcal{T}_p$-basic neighborhood $\supseteq$
    $d_\omega$-ball: Set $\delta = \min_{i \in F}\{\delta_i/(i+1)\}$;
    then $d_\omega(x,y) < \delta$ forces $y_i \in U_i$ for $i \in F$.
\end{itemize}

\textbf{Connectedness:} $(\R^\omega, \mathcal{T}_p)$ is path
connected: $H(t) = ((1-t)x_i + ty_i)_i$ is continuous (coordinatewise
continuous, so continuous in product topology).

\subsection{Box Topology is Disconnected (HW5 Q4)}
\textbf{Key idea:} In $\mathcal{T}_b$, bounded and unbounded
sequences form a nontrivial separation.

$U = \{$bounded sequences$\}$ and $\R^\omega \setminus U$ are both open:
$N_x = \prod_i (x_i - 1, x_i + 1)$ is a box neighborhood that stays
entirely in $U$ (if $x$ bounded) or in the complement (if $x$
unbounded, since $|y_i| \ge |x_i| - 1$).

\subsection{Proving Disconnectedness via Continuous Functions (HW5 Q5)}
\textbf{General technique:} To show $X$ is disconnected, find a
continuous $f : X \to \R$ with $f(X)$ disconnected.

\textbf{Example:} $X = \{(x,0)\} \cup \{(x,1/x) : x > 0\} \subset \R^2$.
Define $f(x,y) = xy$. Then $f(X) = \{0, 1\}$ (disconnected). Hence
$X$ is disconnected.

\subsection{Brouwer Fixed Point Theorem in 1D (HW6 Q3)}
\begin{theorem}
  Every continuous $f : [0,1] \to [0,1]$ has a fixed point.
\end{theorem}
\begin{proof}
  Define $g(x) = f(x) - x$. Then $g(0) = f(0) \ge 0$ and $g(1) = f(1)
  - 1 \le 0$.
  By IVT, $\exists x_0$ with $g(x_0) = 0$, i.e., $f(x_0) = x_0$.
\end{proof}
\textbf{Counterexamples for $(0,1)$ and $(0,1]$:} $f(x) = x/2$ maps
$(0,1) \to (0,1/2) \subset (0,1)$ with no fixed point. Similarly
$f(x) = x/2$ works on $(0,1]$.

\subsection{Connected $+$ Locally Path Connected $\Rightarrow$ Path
Connected (HW6 Q4)}
\begin{definition}
  $X$ is \textbf{locally path connected} if every point has an open
  neighborhood basis of path-connected sets.
\end{definition}

\begin{theorem}
  If $X$ is connected and locally path connected, then $X$ is path connected.
\end{theorem}
\begin{proof}[Proof sketch]
  Fix $x \in X$. Let $S_x = \{y \in X : \exists\text{ path from } x
  \text{ to } y\}$.
  \begin{itemize}
    \item $S_x$ is open: if $y \in S_x$ and $B$ is a path-connected
      neighborhood of $y$, then every $z \in B$ has a path $x \to y
      \to z$, so $B \subseteq S_x$.
    \item $X \setminus S_x$ is open: if $y \notin S_x$ and $B$ is a
      path-connected neighborhood of $y$, then $B \cap S_x =
      \varnothing$ (otherwise concatenate paths to get $y \in S_x$).
  \end{itemize}
  So $S_x$ is clopen. Since $X$ is connected and $S_x \ne
  \varnothing$ (contains $x$), $S_x = X$.
\end{proof}
\textbf{Key technique:} The ``clopen set'' argument. To show a
property holds globally, define $S = \{$points with the property$\}$,
show $S$ is clopen, and invoke connectedness.

\subsection{Non-Homeomorphism via Connectedness (HW6 Q2)}
$(0,\infty)$, $(0,1]$, and $[0,1]$ are pairwise non-homeomorphic.
\textbf{Arguments:}
\begin{itemize}
  \item $[0,1]$ is compact; $(0,\infty)$ and $(0,1]$ are not.
  \item $(0,1] \setminus \{1\} = (0,1)$ is connected, but $[0,1]
    \setminus \{p\}$ is disconnected for $p \in (0,1)$ (interior
    point removal disconnects). Since removing a point from one space
    gives connected and from the other gives disconnected, they
    cannot be homeomorphic.
  \item $(0,\infty) \setminus \{p\}$ is disconnected for all $p$;
    $(0,1] \setminus \{1\} = (0,1)$ is connected.
\end{itemize}
\textbf{General technique:} If $X \cong Y$ and $X \setminus \{p\}
\cong Y \setminus \{f(p)\}$, then topological properties of the
complement (like connectedness) must be preserved.

\subsection{Every Subset Is Compact in Cofinite Topology (HW6 Q5)}
\begin{theorem}
  In $(\R, \mathcal{T}_{\mathrm{cof}})$ where
  $\mathcal{T}_{\mathrm{cof}} = \{\varnothing\} \cup \{X \setminus F
  : F \text{ finite}\}$, every subset is compact.
\end{theorem}
\begin{proof}
  Let $Y \subseteq \R$ and $\{U_i\}$ be an open cover of $Y$. Pick
  any $U_{i_0}$. Then $\R \setminus U_{i_0}$ is finite, so $Y
  \setminus U_{i_0}$ is finite. For each $y \in Y \setminus U_{i_0}$,
  pick $U_{i_y}$ containing $y$. Then $\{U_{i_0}\} \cup \{U_{i_y} : y
  \in Y \setminus U_{i_0}\}$ is a finite subcover.
\end{proof}
\textbf{Key insight:} In the cofinite topology, the complement of any
nonempty open set is finite. So any open cover has at most finitely
many points not covered by a single set.

%=====================================================================
\section{Proof Technique Summary}
%=====================================================================

\begin{enumerate}
  \item \textbf{Double inclusion:} To prove $A = B$, show $A
    \subseteq B$ and $B \subseteq A$.
    Used in: closure $= A \cup LP(A)$, topology equality
    $\mathcal{T}_1 = \mathcal{T}_2$, bijection via injective $+$ surjective.

  \item \textbf{Contraposition:} $P \Rightarrow Q$ $\iff$ $\neg Q
    \Rightarrow \neg P$.
    Used in: proving $f^{-1}$ is not continuous ($f^{-1}$ of open not open).

  \item \textbf{Contradiction:} Assume $\neg$ conclusion, derive contradiction.
    Used in: Cantor's theorem, non-existence of uniform Lipschitz constant.

  \item \textbf{Case analysis:} Break into exhaustive cases.
    Used in: triangle inequality for $d^*(x,y) = \min\{d(x,y),1\}$,
    verifying metric axioms for $d_n = (\ell^1)^n$.

  \item \textbf{Sequence construction:} Build explicit sequences to
    show $x \in \overline{A}$ or $x \in LP(A)$.
    Used in: closure of topologist's sine curve, density in product topology.

  \item \textbf{Free coordinates:} In product topology, unrestricted
    coordinates allow ``witness'' construction.
    Used in: interior/closure computations in $\R^\omega$.

  \item \textbf{Box neighborhood trick:} In box topology, $\prod_i
    (x_i - \varepsilon, x_i + \varepsilon)$ controls every coordinate.
    Used in: bounded/unbounded separation, showing sets are open.

  \item \textbf{Continuous function to $\{0,1\}$:} To show
    disconnectedness, find continuous $f$ with $f(X) = \{0, 1\}$.
    Used in: HW5 Q5 ($f(x,y)=xy$), IVT-based arguments.

  \item \textbf{Topological invariant argument:} To show $X \not\cong
    Y$, find a property preserved by homeomorphism that differs.
    Used in: $[0,1) \not\cong S^1$ (compactness), $(0,1] \not\cong
    [0,1]$ (point-removal connectedness).

  \item \textbf{Clopen set argument:} To show a global property,
    define $S = \{$points with property$\}$, show $S$ is clopen,
    invoke connectedness.
    Used in: connected $+$ locally path connected $\Rightarrow$ path connected.

  \item \textbf{IVT / fixed point:} Define $g(x) = f(x) - x$, check
    signs at endpoints, apply IVT.
    Used in: Brouwer fixed point in 1D.

  \item \textbf{Ball comparison for topology equality:} Show every
    $d_1$-ball contains a $d_2$-ball and vice versa.
    Used in: $d^*$ induces same topology as $d$, metrizability of $\R^\omega$.

  \item \textbf{Monotonicity from additivity:} $A \subseteq B
    \Rightarrow f(A \cup B) = f(A) \cup f(B)$ implies $f(A) \subseteq f(B)$.
    Used in: Kuratowski closure axioms.

  \item \textbf{Characteristic function / indicator:} Encode subsets
    as functions $A \to \{0,1\}$; set operations become pointwise
    arithmetic.
    Used in: bijection $\mathcal{P}(A) \cong 2^A$ (HW1 Q1).

  \item \textbf{Countable union of countable/finite sets:} If
    $S = \bigcup_{n=1}^\infty S_n$ with each $S_n$ countable (or
    finite), then $S$ is countable.
    Used in: algebraic numbers are countable (HW2 Q1), eventually-zero
    sequences are countable (HW2 Q2).

  \item \textbf{Well-definedness on quotient:} When defining a
    function $f$ on equivalence classes $X/{\sim}$, verify $f([x])$
    does not depend on the representative $x$.
    Used in: quotient metric from pseudo-metric (HW2 Q3).

  \item \textbf{Axiom verification (checklist):} Systematically
    verify each axiom: topology (3 axioms), metric (3 axioms),
    equivalence relation (3 axioms), group (4 axioms).
    Used in: intersection of topologies (HW3 Q4), radially open sets
    (HW3 Q5), pseudo-metric quotient (HW2 Q3), symmetric group
    (HW1 Q2).

  \item \textbf{Subbasis method:} To construct the smallest
    topology/smallest equivalence relation containing a family, take
    the intersection of all topologies/relations containing it.
    Used in: topology generated by $\mathcal{T}_1 \cup \mathcal{T}_2$
    (HW3 Q4), generated equivalence relation (HW1 Q4).

  \item \textbf{Counterexample with minimal configuration:} ``Can one
    always\ldots'' questions are often disproved by the smallest
    nontrivial example.
    Used in: no metric on $\mathcal{P}(S)$ extending point distances
    ($|S|=3$, HW2 Q5), $\mathcal{T}_1 \cup \mathcal{T}_2$ not a
    topology ($|S|=3$, HW3 Q4).

  \item \textbf{Induction with auxiliary bound:} When inductively
    bounding a sequence $\{a_n\}$, express $a_{n+1}$ in terms of the
    product $a_1 \cdots a_n$ and use a closed-form bound (geometric
    series).
    Used in: $p_n < 2^{2^n}$ via $p_{n+1} \le p_1 \cdots p_n + 1$
    (HW1 Q3).
\end{enumerate}

%=====================================================================
\section{Additional Proof Techniques in Topology}
%=====================================================================

\subsection{Clopen Set $=$ $\varnothing$ or $X$ in a Connected Space}
\begin{theorem}
  If $(X,\mathcal{T})$ is connected and $A \subseteq X$ is clopen, then
  $A = \varnothing$ or $A = X$.
\end{theorem}
\textbf{Why it matters:} This is the single most versatile tool in
point-set topology. It converts a global statement into a \emph{local}
check (show $A$ is open and closed), then hands the conclusion to
connectedness.

\textbf{Typical proof pattern:}
\begin{enumerate}
  \item Define $A$ to be the set of points with the desired property.
  \item Show $A$ is open (usually easy: use neighborhoods / local structure).
  \item Show $A$ is closed (either: show $X \setminus A$ is open, or
      show $A$ contains all its limit points, or show $A$ is the preimage
    of a closed set).
  \item Check $A \ne \varnothing$ (exhibit one point with the property).
  \item Conclude $A = X$ by connectedness.
\end{enumerate}

\textbf{Applications:}
\begin{itemize}
  \item \textbf{Uniqueness of lifts:} If $f, g : X \to Y$ are continuous
    and agree on a dense set, and $Y$ is Hausdorff, then $A = \{x :
    f(x) = g(x)\}$ is closed (preimage of closed diagonal in $Y \times
    Y$). If also dense, then $A = X$.
  \item \textbf{Identity theorem:} If $f : X \to \R$ is continuous on a
    connected $X$ and $f$ vanishes on an open set, then $f \equiv 0$
    ($\{x : f(x) = 0\}$ is clopen).
  \item \textbf{Locally constant $\Rightarrow$ constant:} If $f : X \to
    Y$ is locally constant (every point has a neighborhood on which $f$
    is constant) and $X$ is connected, then $f$ is constant. For any
    $y_0 \in f(X)$, $f^{-1}(y_0)$ is open (locally constant) and closed
    (preimage of closed $\{y_0\}$ in a Hausdorff $Y$), hence $f^{-1}(y_0) = X$.
  \item \textbf{Connected $+$ locally path connected $\Rightarrow$ path
    connected:} $S_x = \{$points reachable by a path from $x\}$ is clopen.
\end{itemize}

\subsection{Proving a Set $A = \varnothing$}
\begin{itemize}
  \item \textbf{Via connectedness:} Show $A$ is clopen and $A \ne X$.
    If $X$ is connected, the only clopen proper subset is $\varnothing$.
    \textbf{Example:} To show $A \cap B = \varnothing$, show $A$ is
    open, $A \subseteq X \setminus B$, and $B$ is dense (so $X \setminus
    B$ has empty interior).

  \item \textbf{Via interior:} Show $\operatorname{int}(A) = \varnothing$
    and $A$ is open. Then $A = \varnothing$.
    \textbf{Example:} The graph of $f : (0,1] \to \R$ has empty interior
    in $\R^2$ (it's ``too thin'' --- any ball in $\R^2$ contains points
    not on the graph). Hence the graph is nowhere dense.

  \item \textbf{Via closure:} Show $\overline{A} = \varnothing$. Since
    $A \subseteq \overline{A}$, this forces $A = \varnothing$.

  \item \textbf{By contradiction:} Assume $x \in A$, derive a
    contradiction with a known property.
    \textbf{Example:} $\operatorname{int}(\overline{A}) = \varnothing$
    when $A = \Z \subset \R$: any open interval $(a,b)$ contains
    irrationals $\notin \overline{\Z} = \Z$.
\end{itemize}

\subsection{Proving a Set $A = X$ (Global Property)}
\begin{itemize}
  \item \textbf{Clopen $+$ nonempty:} Show $A$ is clopen and $A \ne
    \varnothing$ in a connected $X$.
  \item \textbf{Dense $+$ open:} If $A$ is dense ($\overline{A} = X$)
    and open, then $A = X$. \textbf{Reason:} $\overline{A} = X$ and $A$
    is open $\Rightarrow$ $X \setminus A$ is closed. If $X \setminus A
    \ne \varnothing$, then $A \subseteq X \setminus (X \setminus A)
    \subsetneq X$, contradicting density.
  \item \textbf{Via complement $= \varnothing$:} Show $X \setminus A =
    \varnothing$ using the methods above.
\end{itemize}

\subsection{Preimage / Inverse Image Arguments}
The preimage $f^{-1}(U)$ is extremely well-behaved:
\[
  f^{-1}\!\Bigl(\bigcup_\alpha U_\alpha\Bigr) = \bigcup_\alpha
  f^{-1}(U_\alpha), \qquad
  f^{-1}\!\Bigl(\bigcap_\alpha U_\alpha\Bigr) = \bigcap_\alpha
  f^{-1}(U_\alpha), \qquad
  f^{-1}(Y \setminus U) = X \setminus f^{-1}(U).
\]
These hold for \emph{any} function $f$ (no continuity needed).
Continuity gives the additional conclusion: ``$U$ open $\Rightarrow$
$f^{-1}(U)$ open'' (and dually for closed).

\textbf{Common patterns:}
\begin{itemize}
  \item \textbf{Level sets:} $f^{-1}(\{c\})$ is closed if $f$ is
    continuous and $\{c\}$ is closed (e.g.\ $\R$ is Hausdorff).
    Used to prove that the zero set of a continuous function is closed.
  \item \textbf{Preimage of a separation:} If $Y = U \cup V$ is a
    separation, then $f^{-1}(U) \cup f^{-1}(V) = X$ is a separation
    (if $f$ is continuous and surjective). Used to prove: continuous
    surjective image of connected is connected (contrapositive).
  \item \textbf{Preimage of Hausdorff:} If $Y$ is Hausdorff, then the
    \textbf{diagonal} $\Delta = \{(y,y) : y \in Y\}$ is closed in $Y
    \times Y$. If $f,g : X \to Y$ are continuous, then $\{x : f(x) =
    g(x)\} = (f \times g)^{-1}(\Delta)$ is closed.
\end{itemize}

\subsection{Open Set Description (``Pick a Point and Find a Ball'')}
To show a set $U$ is open in a metric space:
\begin{enumerate}
  \item Take an arbitrary $x \in U$.
  \item Find $\varepsilon > 0$ such that $B(x, \varepsilon) \subseteq U$.
  \item Conclude $U$ is open.
\end{enumerate}
\textbf{This is the definition}, but the art lies in choosing
$\varepsilon$.

\textbf{Typical choices of $\varepsilon$:}
\begin{itemize}
  \item $\varepsilon = r - d(x, x_0)$ to show $B(x_0, r)$ is open.
  \item $\varepsilon = \min\{\varepsilon_1, \ldots, \varepsilon_k\}$
    for finite intersections of open sets.
  \item $\varepsilon$ from the bounded metric: $d^*(x,y) < \varepsilon$
    with $\varepsilon \le 1$ so that $d^* = d$.
\end{itemize}

For a general topological space, replace ``find a ball'' with ``find a
basic open set $B \in \mathcal{B}$ with $x \in B \subseteq U$.''

\subsection{Neighborhood Basis / Local Basis Argument}
A \textbf{neighborhood basis} at $x$ is a collection $\{N_\alpha\}$ of
neighborhoods of $x$ such that every neighborhood of $x$ contains some
$N_\alpha$.

\textbf{Use:} To verify a local property at $x$ (continuity, openness,
limit point, etc.), it suffices to check on a neighborhood basis.

\textbf{Examples:}
\begin{itemize}
  \item In a metric space: $\{B(x, 1/n) : n \in \N\}$ is a countable
    neighborhood basis.
  \item In the product topology on $\prod X_i$: basic neighborhoods of
    $x$ are $\prod U_i$ where $U_i = X_i$ for all but finitely many $i$.
  \item To show $x \in \overline{A}$: check that every basic
    neighborhood of $x$ meets $A$.
\end{itemize}

\subsection{Closure Characterization Switching}
The closure has many equivalent descriptions. Pick the one that makes
your proof easiest:
\begin{enumerate}
  \item $\overline{A} = \bigcap\{C \supseteq A : C \text{ closed}\}$
    (intersection definition --- useful for proving $\overline{A}$ is closed).
  \item $x \in \overline{A} \iff$ every open neighborhood of $x$ meets $A$
    (neighborhood definition --- useful for showing $x \in \overline{A}$).
  \item $x \in \overline{A} \iff$ $\exists$ sequence $(a_n) \subseteq A$
    with $a_n \to x$ (sequential characterization --- only valid in
    first-countable spaces, e.g.\ metric spaces).
  \item $\overline{A} = A \cup LP(A) = A \cup A'$ (via limit points).
  \item $\overline{A} = A \cup \partial A$ (via boundary).
\end{enumerate}

\textbf{Strategy:}
\begin{itemize}
  \item To show $x \in \overline{A}$: use (2) or (3) (construct a
    sequence or argue every neighborhood meets $A$).
  \item To show $A$ is dense ($\overline{A} = X$): use (2) --- show
    every nonempty open set meets $A$. E.g.\ $\Q$ is dense in $\R$:
    every interval contains a rational.
  \item To show $\overline{A} \subseteq B$: use (1) --- show $B$ is a
    closed superset of $A$, hence $\overline{A} \subseteq B$ (since
    $\overline{A}$ is the \emph{smallest} closed superset).
  \item To show $\overline{A} = \overline{B}$: use double inclusion with
    the neighborhood criterion. E.g.\ $A$ dense in $B$ and $B$ dense in
    $X$ $\Rightarrow$ $\overline{A} = \overline{B} = X$.
\end{itemize}

\subsection{``It Suffices to Show'' --- Reducing to a Simpler Claim}
Many topology proofs use the structure ``to prove $P$, it suffices to
show $Q$,'' where $Q$ is a simpler, more local statement.

\textbf{Common reductions:}
\begin{itemize}
  \item \textbf{Continuity:} To prove $f$ is continuous globally, it
    suffices to check $f^{-1}(B)$ is open for every subbasis element $B$
    (instead of every open set).
  \item \textbf{Topology equality $\mathcal{T}_1 = \mathcal{T}_2$:} It
    suffices to show every $\mathcal{T}_1$-basic open set is
    $\mathcal{T}_2$-open and vice versa.
  \item \textbf{Density:} To show $A$ is dense, it suffices to show $A$
    meets every basis element (not every open set).
  \item \textbf{Nowhere dense:} To show $\operatorname{int}(\overline{A})
    = \varnothing$, it suffices to show every nonempty open set contains
    a nonempty open subset disjoint from $A$.
\end{itemize}

\subsection{Contrapositive in Topology}
Many topological statements are most naturally proved in contrapositive form.

\textbf{Pattern:} ``$X$ connected $\Rightarrow$ $f(X)$ connected'' is
proved as: ``$f(X)$ disconnected $\Rightarrow$ $X$ disconnected.'' If
$f(X) = U \cup V$ is a separation, then $f^{-1}(U) \cup f^{-1}(V) = X$
is a separation of $X$ (preimages preserve openness, disjointness, and
covering).

\textbf{Other examples:}
\begin{itemize}
  \item ``Continuous image of compact is compact'' $\equiv$
    ``$f(X)$ not compact $\Rightarrow$ $X$ not compact.''
  \item ``$f$ continuous $\Rightarrow$ $f^{-1}$(closed) is closed''
    $\equiv$ ``$f^{-1}(F)$ not closed $\Rightarrow$ $f$ not continuous.''
  \item ``Homeomorphism preserves topological invariants'' $\equiv$
    ``$X, Y$ differ on some invariant $\Rightarrow$ $X \not\cong Y$.''
\end{itemize}

\subsection{Compactness Arguments}

\begin{theorem}[Continuous image of compact is compact]
  If $f : X \to Y$ is continuous and $X$ is compact, then $f(X)$ is compact.
\end{theorem}
\begin{proof}
  Let $\{V_\alpha\}$ be an open cover of $f(X)$. Since $f$ is continuous,
  each $f^{-1}(V_\alpha)$ is open in $X$, and $\{f^{-1}(V_\alpha)\}$ covers $X$
  (for each $x \in X$, $f(x) \in V_\alpha$ for some $\alpha$, so $x
  \in f^{-1}(V_\alpha)$).
  Since $X$ is compact, extract a finite subcover $X =
  f^{-1}(V_{\alpha_1}) \cup \cdots \cup f^{-1}(V_{\alpha_k})$.
  Then $f(X) \subseteq V_{\alpha_1} \cup \cdots \cup V_{\alpha_k}$, a
  finite subcover.
\end{proof}

\begin{theorem}[Closed subset of compact is compact]
  If $K$ is compact and $F \subseteq K$ is closed, then $F$ is compact.
\end{theorem}
\begin{proof}
  Let $\{U_\alpha\}$ be an open cover of $F$. Then $\{U_\alpha\} \cup
  \{K \setminus F\}$
  is an open cover of $K$ (since $F$ is closed, $K \setminus F$ is open).
  By compactness of $K$, extract a finite subcover.
  Removing $K \setminus F$ from this finite list (if present) gives a finite
  subcover of $F$ from $\{U_\alpha\}$.
\end{proof}

\begin{theorem}[Compact subset of Hausdorff space is closed]
  In a Hausdorff space $Y$, every compact subset $K \subseteq Y$ is closed.
\end{theorem}
\begin{proof}
  We show $Y \setminus K$ is open. Fix $y \notin K$.
  For each $x \in K$, by Hausdorff, pick disjoint open $U_x \ni x$
  and $V_x \ni y$.
  Then $\{U_x\}_{x \in K}$ covers $K$. Since $K$ is compact, extract
  a finite subcover:
  $K \subseteq U_{x_1} \cup \cdots \cup U_{x_n}$.
  Set $V = V_{x_1} \cap \cdots \cap V_{x_n}$. Then $V$ is an open
  neighborhood of $y$,
  and $V \cap U_{x_i} = \varnothing$ for each $i$, so $V \cap K = \varnothing$.
  Hence $y \in V \subseteq Y \setminus K$, proving $Y \setminus K$ is open.
\end{proof}

\begin{theorem}[Tube Lemma]
  Let $X \times Y$ have the product topology with $Y$ compact.
  If $N$ is an open set containing the ``slice'' $\{x_0\} \times Y$,
  then there exists an open neighborhood $W$ of $x_0$ such that
  $W \times Y \subseteq N$ (a ``tube'' around the slice).
\end{theorem}
\begin{proof}
  For each $y \in Y$, $(x_0, y) \in N$. Since $N$ is open, there is a
  basic open set
  $U_y \times V_y \subseteq N$ with $x_0 \in U_y$, $y \in V_y$.
  Then $\{V_y\}_{y \in Y}$ covers $Y$. By compactness, extract a finite subcover
  $Y = V_{y_1} \cup \cdots \cup V_{y_n}$. Set $W = U_{y_1} \cap
  \cdots \cap U_{y_n}$.
  For any $(x, y) \in W \times Y$: $y \in V_{y_j}$ for some $j$, and
  $x \in W \subseteq U_{y_j}$, so $(x, y) \in U_{y_j} \times V_{y_j}
  \subseteq N$.
\end{proof}

\begin{theorem}[Heine--Borel]
  A subset $A \subseteq \R^n$ is compact (in the standard topology)
  if and only if $A$ is closed and bounded.
\end{theorem}
\begin{proof}
  ($\Rightarrow$) Since $\R^n$ is Hausdorff, compact $A$ is closed.
  Since $A \subseteq \bigcup_{n=1}^\infty B(0, n)$ is an open cover,
  compactness gives a finite subcover, so $A \subseteq B(0, N)$ for
  some $N$, i.e.\ bounded.

  ($\Leftarrow$) If $A$ is closed and bounded, then $A \subseteq [-M,
  M]^n$ for some $M$.
  By Tychonoff (finite case), $[-M, M]^n$ is compact (product of
  compact spaces).
  Since $A$ is a closed subset of the compact $[-M, M]^n$, $A$ is compact.
  The base case $[a,b]$ is compact: given open cover $\{U_\alpha\}$,
  set $c = \sup\{t \in [a,b] : [a,t]$ has a finite subcover$\}$.
  Then $c \in U_\beta$ for some $\beta$, so $[c-\varepsilon,
  c+\varepsilon] \subseteq U_\beta$,
  extending the finite subcover past $c$ unless $c = b$. Hence $c =
  b$ and $[a,b]$ is compact.
\end{proof}

\textbf{Finite subcover trick:} Given an open cover of a compact
space, extract a finite subcover and use finiteness (e.g.\ take
minimum of finitely many positive numbers, which is still positive).
\textbf{Example:} Heine--Cantor theorem: continuous on compact
$\Rightarrow$ uniformly continuous. Cover $K$ by $B(x,
\delta_x/2)$; extract finite subcover; set $\delta =
\min\{\delta_{x_i}/2\} > 0$.

\textbf{Vacuous truth (空真):} ``$\forall x \in \varnothing, P(x)$'' is
true (nothing to check). This is why $\varnothing$ is open.

\textbf{For almost all (除有限多个外全部):} In the product topology, $U_i =
X_i$ for all but finitely many $i$.

%=====================================================================
\newpage
\section{Terminology: Chinese--English Glossary (中英对照术语表)}
%=====================================================================

\subsection*{集合论与公理 (Set Theory and Axioms)}

\begin{tabular}{@{}lp{7.5cm}c@{}}
  \hline
  \textbf{English} & \textbf{中文} & \textbf{Ref.} \\
  \hline
  Axiom of Extensionality          & 外延公理 & \S\ref{def:extensionality} \\
  Axiom of Empty Set               & 空集公理 & \S\ref{def:empty-set} \\
  Axiom of Pairing                 & 配对公理 & \S\ref{def:pairing} \\
  Axiom of Union                   & 并集公理 & \S\ref{def:union-axiom} \\
  Axiom of Power Set               & 幂集公理 & \S\ref{def:power-set-axiom} \\
  Axiom of Infinity                & 无穷公理 & \S\ref{def:infinity} \\
  Axiom Schema of Separation       & 分离公理（概括公理）& \S\ref{def:separation} \\
  Axiom Schema of Replacement      & 替换公理 & \S\ref{def:replacement} \\
  Axiom of Regularity (Foundation) & 正则公理（基础公理）& \S\ref{def:regularity} \\
  Axiom of Choice (AC)             & 选择公理 & \S\ref{def:ac} \\
  Russell's paradox                & 罗素悖论 & \S\ref{def:russell} \\
  Forcing                          & 力迫法 & \S\ref{def:forcing} \\
  Power set $\mathcal{P}(A)$       & 幂集 & \S\ref{def:power-set-axiom} \\
  Characteristic / indicator function $\chi_A$ & 特征函数 / 指示函数 &
  \S\ref{def:power-set-axiom} \\
  Equivalence relation             & 等价关系 & --- \\
  Equivalence class                & 等价类 & --- \\
  Transitive closure               & 传递闭包 & --- \\
  Symmetric group $S_n$            & 对称群 & --- \\
  Bijection                        & 双射 / 一一对应 & --- \\
  Injection                        & 单射 & --- \\
  Surjection                       & 满射 & --- \\
  \hline
\end{tabular}

\subsection*{可数性与基数 (Countability and Cardinality)}

\begin{tabular}{@{}lp{7.5cm}c@{}}
  \hline
  \textbf{English} & \textbf{中文} & \textbf{Ref.} \\
  \hline
  Finite set                       & 有限集 & \S\ref{def:finite} \\
  Infinite set                     & 无限集 & \S\ref{def:infinite} \\
  Countably infinite               & 可数无限 & \S\ref{def:countably-infinite} \\
  Countable                        & 可数 & \S\ref{def:countable} \\
  Uncountable                      & 不可数 & \S\ref{def:uncountable} \\
  Cantor diagonal argument         & Cantor 对角线论证 & \S\ref{def:cantor-diag} \\
  Cantor's theorem                 & Cantor
  定理（$|A|<|\mathcal{P}(A)|$）& \S\ref{def:cantor-thm} \\
  Cardinal number                  & 基数 & \S\ref{def:cardinality} \\
  Cardinality                      & 势 / 基数 & \S\ref{def:cardinality} \\
  Same cardinality (equipotent)    & 等势 & \S\ref{def:cardinality} \\
  $\aleph_0$ (aleph-nought)        & 可数无限集的基数 & \S\ref{def:aleph} \\
  $\aleph_1$                       & 连续统的基数 & \S\ref{def:aleph} \\
  Continuum Hypothesis (CH)        & 连续统假设 & \S\ref{def:ch} \\
  Algebraic number                 & 代数数 & --- \\
  Transcendental number            & 超越数 & --- \\
  \hline
\end{tabular}

\subsection*{拓扑空间 (Topological Spaces)}

\begin{tabular}{@{}lp{7.5cm}c@{}}
  \hline
  \textbf{English} & \textbf{中文} & \textbf{Ref.} \\
  \hline
  Topological space                & 拓扑空间 & \S\ref{def:topology} \\
  Topology                         & 拓扑 & \S\ref{def:topology} \\
  Open set                         & 开集 & \S\ref{def:open} \\
  Closed set                       & 闭集 & \S\ref{def:closed} \\
  Clopen set                       & 既开又闭集 & \S\ref{def:clopen} \\
  Trivial (indiscrete) topology    & 平凡拓扑（密着拓扑）& \S\ref{def:trivial} \\
  Discrete topology                & 离散拓扑 & \S\ref{def:discrete-topo} \\
  Cofinite topology                & 余有限拓扑 & \S\ref{def:cofinite} \\
  Semicontinuous topology          & 半无穷拓扑 & \S\ref{def:semicontinuous} \\
  Finer / larger topology          & 更细 / 更大的拓扑 & \S\ref{def:finer-coarser} \\
  Coarser / smaller topology       & 更粗 / 更小的拓扑 & \S\ref{def:finer-coarser} \\
  Generated topology               & 生成的拓扑 & \S\ref{def:generate-topo} \\
  Radially open                    & 径向开集 & --- \\
  \hline
\end{tabular}

\subsection*{基 (Basis)}

\begin{tabular}{@{}lp{7.5cm}c@{}}
  \hline
  \textbf{English} & \textbf{中文} & \textbf{Ref.} \\
  \hline
  Basis (base)                     & 基 & \S\ref{def:basis} \\
  Subbasis (subbase)               & 子基 & \S\ref{def:subbasis} \\
  Generate a topology              & 生成拓扑 & \S\ref{def:generate-topo} \\
  Countable basis (second-countable) & 可数基（第二可数）&
  \S\ref{def:second-countable} \\
  Open rectangle / open ball       & 开矩形 / 开球 & \S\ref{def:open-rect} \\
  \hline
\end{tabular}

\subsection*{度量空间 (Metric Spaces)}

\begin{tabular}{@{}lp{7.5cm}c@{}}
  \hline
  \textbf{English} & \textbf{中文} & \textbf{Ref.} \\
  \hline
  Metric space                     & 度量空间 & \S\ref{def:metric} \\
  Metric / distance function       & 度量 / 距离函数 & \S\ref{def:metric} \\
  Discrete metric                  & 离散度量 & \S\ref{def:discrete-metric} \\
  Open ball                        & 开球 & \S\ref{def:balls} \\
  Closed ball                      & 闭球 & \S\ref{def:balls} \\
  Bounded set                      & 有界集 & \S\ref{def:bounded} \\
  Bounded metric                   & 有界度量 & \S\ref{def:bounded-metric} \\
  Equivalent metrics               & 等价度量 & \S\ref{def:equiv-metrics} \\
  Metrizable (topology)            & 可度量化的（拓扑）& \S\ref{def:metrizable} \\
  Euclidean metric                 & 欧几里得度量 & \S\ref{def:euclidean} \\
  Taxicab metric ($\ell^1$)        & 曼哈顿度量 / 出租车度量 & \S\ref{def:taxicab} \\
  Supremum metric ($\ell^\infty$)  & 最大值度量 & \S\ref{def:supremum} \\
  Norm                             & 范数 & \S\ref{def:norm} \\
  Pseudo-metric                    & 伪度量 & --- \\
  Quotient metric / quotient space & 商度量 / 商空间 & --- \\
  Well-defined                     & 良定义 & --- \\
  \hline
\end{tabular}

\subsection*{分离公理 (Separation Axioms)}

\begin{tabular}{@{}lp{7.5cm}c@{}}
  \hline
  \textbf{English} & \textbf{中文} & \textbf{Ref.} \\
  \hline
  $T_0$ / Kolmogorov space         & Kolmogorov 空间 & \S\ref{def:T0} \\
  $T_1$ / Fr\'echet space          & Fr\'echet 空间 & \S\ref{def:T1} \\
  $T_2$ / Hausdorff space          & Hausdorff 空间 & \S\ref{def:T2} \\
  $T_3$ / Regular space            & 正则空间 & \S\ref{def:T3} \\
  $T_{3\frac{1}{2}}$ / Tychonoff (completely regular) & 完全正则空间 &
  \S\ref{def:T3h} \\
  $T_4$ / Normal space             & 正规空间 & \S\ref{def:T4} \\
  Urysohn's Lemma                  & Urysohn 引理 & \S\ref{def:T4} \\
  Tietze Extension Theorem         & Tietze 扩张定理 & \S\ref{def:T4} \\
  Urysohn Metrization Theorem      & Urysohn 可度量化定理 &
  \S\ref{thm:urysohn-metrization} \\
  \hline
\end{tabular}

\subsection*{贝尔纲定理 (Baire Category Theorem)}

\begin{tabular}{@{}lp{7.5cm}c@{}}
  \hline
  \textbf{English} & \textbf{中文} & \textbf{Ref.} \\
  \hline
  Nowhere dense set                & 疏朗集 / 无处稠密集 &
  \S\ref{def:nowhere-dense-baire} \\
  Meager / first category set      & 第一纲集 & \S\ref{def:meager} \\
  Second category set              & 第二纲集 & \S\ref{def:non-meager} \\
  Baire space                      & Baire 空间 & \S\ref{def:baire-space} \\
  Baire Category Theorem           & 贝尔纲定理 & \S\ref{thm:baire} \\
  $G_\delta$ set                   & $G_\delta$ 集 &
  \S\ref{def:baire-applications} \\
  \hline
\end{tabular}

\subsection*{子空间 (Subspace)}

\begin{tabular}{@{}lp{7.5cm}c@{}}
  \hline
  \textbf{English} & \textbf{中文} & \textbf{Ref.} \\
  \hline
  Subspace topology                & 子空间拓扑 & \S\ref{def:subspace-topo} \\
  Subspace                         & 子空间 & \S\ref{def:subspace} \\
  Hereditary property              & 遗传性质 & \S\ref{def:hereditary} \\
  \hline
\end{tabular}

\subsection*{闭包、内部与边界 (Closure, Interior, Boundary)}

\begin{tabular}{@{}lp{7.5cm}c@{}}
  \hline
  \textbf{English} & \textbf{中文} & \textbf{Ref.} \\
  \hline
  Closure                          & 闭包 & \S\ref{def:closure} \\
  Interior                         & 内部 & \S\ref{def:interior} \\
  Boundary                         & 边界 & \S\ref{def:boundary} \\
  Neighborhood                     & 邻域 & \S\ref{def:neighborhood} \\
  Limit point (cluster/accumulation point) & 极限点（聚点/附着点）&
  \S\ref{def:limit-point} \\
  Isolated point                   & 孤立点 & \S\ref{def:isolated} \\
  Derived set                      & 导集 & \S\ref{def:derived} \\
  Dense set                        & 稠密集 & \S\ref{def:dense} \\
  Nowhere dense                    & 疏集 / 无处稠密 & \S\ref{def:nowhere-dense} \\
  Regularly open                   & 正则开集 & \S\ref{def:regularly-open} \\
  \hline
\end{tabular}

\subsection*{连续映射 (Continuous Maps)}

\begin{tabular}{@{}lp{7.5cm}c@{}}
  \hline
  \textbf{English} & \textbf{中文} & \textbf{Ref.} \\
  \hline
  Continuous function / map        & 连续函数 / 连续映射 & \S\ref{def:continuous} \\
  Homeomorphism                   & 同胚 & \S\ref{def:homeomorphism} \\
  Topological invariant           & 拓扑不变量 & \S\ref{def:topo-invariant} \\
  Embedding                       & 嵌入 & \S\ref{def:embedding} \\
  Pasting (gluing) lemma          & 粘接引理 & \S\ref{def:pasting} \\
  Restriction                     & 限制 & \S\ref{def:restriction} \\
  Composition                     & 复合 & \S\ref{def:composition} \\
  Coordinate function             & 坐标函数 & \S\ref{def:coordinate} \\
  \hline
\end{tabular}

\subsection*{积拓扑与箱拓扑 (Product and Box Topology)}

\begin{tabular}{@{}lp{7.5cm}c@{}}
  \hline
  \textbf{English} & \textbf{中文} & \textbf{Ref.} \\
  \hline
  Product topology                 & 积拓扑 & \S\ref{def:product-topo} \\
  Box topology                    & 箱拓扑 & \S\ref{def:box-topo} \\
  Projection map                  & 投影映射 & \S\ref{def:projection} \\
  Cylinder (cylinder set)         & 柱集 & \S\ref{def:cylinder} \\
  Universal property              & 泛性质 & \S\ref{def:universal} \\
  Eventually-zero sequence        & 终零序列 & \S\ref{def:eventually-zero} \\
  Bounded sequence                & 有界序列 & \S\ref{def:bounded-seq} \\
  Torus ($S^1 \times S^1$)        & 环面 & \S\ref{def:torus} \\
  \hline
\end{tabular}

\subsection*{连通性 (Connectedness)}

\begin{tabular}{@{}lp{7.5cm}c@{}}
  \hline
  \textbf{English} & \textbf{中文} & \textbf{Ref.} \\
  \hline
  Connected space                  & 连通空间 & \S\ref{def:connected} \\
  Disconnected space              & 不连通空间 & \S\ref{def:disconnected} \\
  Totally disconnected            & 完全不连通 & \S\ref{def:totally-disc} \\
  Separation                      & 分离 & \S\ref{def:separation-conn} \\
  Connected component             & 连通分支 & \S\ref{def:component} \\
  Path connected                  & 道路连通 & \S\ref{def:path-connected} \\
  Path component                  & 道路连通分支 & \S\ref{def:path-component} \\
  Path / curve                    & 道路 / 曲线 & \S\ref{def:path} \\
  Topologist's sine curve         & 拓扑学家的正弦曲线 & \S\ref{def:sine-curve} \\
  Intermediate Value Theorem (IVT)& 介值定理 & \S\ref{def:ivt} \\
  \hline
\end{tabular}

\subsection*{紧致性 (Compactness)}

\begin{tabular}{@{}lp{7.5cm}c@{}}
  \hline
  \textbf{English} & \textbf{中文} & \textbf{Ref.} \\
  \hline
  Compact space                   & 紧致空间 / 紧空间 & \S\ref{def:compact} \\
  Open cover                      & 开覆盖 & \S\ref{def:open-cover} \\
  Finite subcover                 & 有限子覆盖 & \S\ref{def:finite-subcover} \\
  Subcover                        & 子覆盖 & \S\ref{def:subcover} \\
  Heine--Borel theorem            & Heine--Borel 定理 & \S\ref{def:heine-borel} \\
  \hline
\end{tabular}

\subsection*{常用短语与表达 (Common Phrases)}

\begin{tabular}{@{}lp{7.5cm}c@{}}
  \hline
  \textbf{English} & \textbf{中文} & \textbf{Ref.} \\
  \hline
  Endowed with / equipped with    & 赋予……拓扑 & --- \\
  With respect to                 & 关于…… & --- \\
  Induced topology                & 诱导拓扑 & --- \\
  Show that / Prove that          & 证明 / 说明 & --- \\
  Determine whether               & 判断是否 & --- \\
  Find / Compute                  & 求出 / 计算 & --- \\
  Justify your answer             & 证明你的结论 & --- \\
  Is it possible to find          & 是否存在 & --- \\
  Conclude that                   & 得出结论 & --- \\
  By contradiction                & 反证法 & --- \\
  Contraposition                  & 逆否命题 & --- \\
  Without loss of generality      & 不失一般性 & --- \\
  If and only if (iff)            & 当且仅当 & --- \\
  Suppose / Assume                & 假设 & --- \\
  Hence / Therefore               & 因此 & --- \\
  It follows that                 & 由此可得 & --- \\
  It suffices to show             & 只需证明 & --- \\
  Note that                       & 注意到 & --- \\
  Define                          & 定义 & --- \\
  Denote by                       & 记……为 & --- \\
  For almost all                  & 除有限多个外全部 & --- \\
  Vacuous truth                   & 空真 & --- \\
  \hline
\end{tabular}

\end{document}
